% Main file for the skripta — includes chapters 1..8
\documentclass[12pt,a4paper]{report}

\usepackage[utf8]{inputenc}
\usepackage[T1]{fontenc}
\usepackage[croatian]{babel}
\usepackage{amsmath, amssymb, amsthm}
\usepackage{geometry}
\usepackage{hyperref}
\usepackage{enumitem}
\usepackage{xcolor}
\usepackage{tcolorbox}
\usepackage{parskip}
\usepackage{ifthen}
\tcbuselibrary{theorems, skins, breakable}

\geometry{margin=2.8cm, top=3cm, bottom=3cm}

\hypersetup{
  colorlinks=true,
  linkcolor=blue!60!black,
  urlcolor=blue!60!black
}

% -------------------------------------------------------
%  Boje, okruzenja i makroi (iste kao u poglavljima)
% -------------------------------------------------------
\definecolor{defbg}{RGB}{235, 244, 255}
\definecolor{defborder}{RGB}{100, 160, 220}
\definecolor{exbg}{RGB}{240, 252, 240}
\definecolor{exborder}{RGB}{100, 180, 100}
\definecolor{keybg}{RGB}{255, 248, 225}
\definecolor{keyborder}{RGB}{210, 170, 50}
\definecolor{taskbg}{RGB}{250, 240, 255}
\definecolor{taskborder}{RGB}{160, 100, 200}

\newtcolorbox[auto counter, number within=section]{definicija}[1][]{
  colback=defbg, colframe=defborder,
  fonttitle=\bfseries\sffamily,
  title={Definicija~\thetcbcounter \ifthenelse{\equal{#1}{}}{}{~---~#1}},
  breakable, enhanced, left=6pt, right=6pt, top=4pt, bottom=4pt,
  arc=4pt
}

\newtcolorbox[auto counter, number within=section]{primjer}[1][]{
  colback=exbg, colframe=exborder,
  fonttitle=\bfseries\sffamily,
  title={Primjer~\thetcbcounter \ifthenelse{\equal{#1}{}}{}{~---~#1}},
  breakable, enhanced, left=6pt, right=6pt, top=4pt, bottom=4pt,
  arc=4pt
}

\newtcolorbox{kljucno}[1][Ključna ideja]{
  colback=keybg, colframe=keyborder,
  fonttitle=\bfseries\sffamily, title=#1,
  breakable, enhanced, left=6pt, right=6pt, top=4pt, bottom=4pt,
  arc=4pt
}

\newtcolorbox[auto counter, number within=section]{zadatak}[1][]{
  colback=taskbg, colframe=taskborder,
  fonttitle=\bfseries\sffamily,
  title={Zadatak~\thetcbcounter \ifthenelse{\equal{#1}{}}{}{~---~#1}},
  breakable, enhanced, left=6pt, right=6pt, top=4pt, bottom=4pt,
  arc=4pt
}

\newtheorem*{napomena}{Napomena}

% theorem-like environments used in chapters
\newtheorem{propozicija}{Propozicija}[section]

\newcommand{\Sig}{\Sigma}
\newcommand{\eps}{\varepsilon}

\usepackage{titlesec}
\titleformat{\section}{\large\bfseries\sffamily}{}{0em}{\thesection\quad}[\vspace{-0.3em}\rule{\linewidth}{0.4pt}]
\titleformat{\subsection}{\normalsize\bfseries\sffamily}{}{0em}{\thesubsection\quad}

\begin{document}

% =========================================================
%  Naslovna stranica
% =========================================================
\begin{titlepage}
\thispagestyle{empty}
\centering

\vspace*{0.8cm}

% --- Vrh: institucija ---
{\sffamily\large Sveučilište Jurja Dobrile u Puli}\\[0.4em]
{\sffamily\normalsize Fakultet informatike}\\[0.4em]
{\sffamily\small Preddiplomski studij informatike}

\vspace{2.8cm}

% --- Dekorativna gornja crta ---
{\color{defborder}\rule{0.85\linewidth}{2pt}}\\[3pt]
{\color{defborder}\rule{0.85\linewidth}{0.4pt}}

\vspace{1.2cm}

% --- Glavni naslov ---
{\fontsize{46}{52}\selectfont\bfseries\sffamily%
  \color{defborder!75!black} Formalni jezici\par}

\vspace{0.5em}

{\Large\sffamily\itshape skripta za pripremu ispita}

\vspace{1.2cm}

% --- Dekorativna donja crta (zrcalno) ---
{\color{defborder}\rule{0.85\linewidth}{0.4pt}}\\[3pt]
{\color{defborder}\rule{0.85\linewidth}{2pt}}

\vspace{2.2cm}

% --- Moto / kratki citat ---
\begin{minipage}{0.78\textwidth}
\centering
{\Large\itshape ``Pripada li ova riječ ovome jeziku?''}\\[0.8em]
{\sffamily\normalsize Jedno pitanje. Pet klasa jezika.\\
Beskonačno mnogo posljedica.}
\end{minipage}

\vfill

% --- Autor i godina, u boksu ---
\begin{tcolorbox}[
  width=0.62\textwidth,
  colback=defbg,
  colframe=defborder,
  arc=5pt,
  boxrule=0.6pt,
  halign=center,
  left=10pt, right=10pt, top=8pt, bottom=8pt
]
{\large\bfseries\sffamily Petar Prenc}\\[0.3em]
{\sffamily\small Akademska godina 2025./2026.}
\end{tcolorbox}

\vspace{0.8cm}

{\sffamily\small Pula, \today}

\vspace{0.5cm}
\end{titlepage}

% =========================================================
%  Predgovor
% =========================================================
\chapter*{Predgovor}
\addcontentsline{toc}{chapter}{Predgovor}

\begin{quote}
\itshape
Najbolji način da nešto stvarno naučiš jest da to objasniš nekom drugom.
Ova skripta upravo je to: pokušaj da samom sebi --- a sad i vama --- na
jednom mjestu objasnim što su to formalni jezici, na način koji se može
pročitati u jednom dahu, bez otvaranja desetak izvora paralelno.
\end{quote}


% ================================================================
\section*{Što je ovo i zašto postoji}
% ================================================================

Skripta je nastala iz vrlo obične potrebe. Pripremajući se za ispit iz
kolegija \emph{Formalni jezici}, primijetio sam da je gradivo na
predavanjima i u udžbeniku raspršeno: definicije ovdje, primjeri ondje,
dokazi negdje između. Da bih ga svladao, morao sam ga prvo okupiti na
jednom mjestu --- vlastitim riječima, slijedeći jednu jasnu nit od
početka do kraja. Rezultat te vježbe je tekst pred vama.

Cilj je skroman: \textbf{biti razumljiv prijateljski vodič kroz gradivo},
ništa više. Ne pretendira zamijeniti službeni udžbenik niti biti
referenca posljednje riječi u svakom dokazu. Želi biti ono što bih i sam
htio imati u rukama tjedan dana prije ispita --- kratak, jasan,
pun primjera i s naglaskom na intuiciji.

% ================================================================
\section*{Što skripta sadrži}
% ================================================================

Materijal prati \textbf{Chomskyjevu hijerarhiju} odozdo prema gore.
Krećemo od najjednostavnijih objekata --- riječi i jezika kao skupova
riječi --- i postupno gradimo sve moćnije formalizme. Svako poglavlje
prirodno se nastavlja na prethodno: regularni izrazi vode do automata,
automati i izrazi pokazuju se ekvivalentnima, ekvivalencija se proteže
i na linearne gramatike, a kad ti formalizmi postanu preuski, prelazimo
na bezkontekstne, kontekstne i konačno Turingove strojeve.

\begin{kljucno}[Plan skripte]
\begin{itemize}[itemsep=2pt]
  \item \textbf{Poglavlje 1} --- Jezici i osnovne operacije: abeceda,
        riječ, $\Sig^*$, konkatenacija, Kleeneva zvijezda, Chomskyjeva
        hijerarhija.
  \item \textbf{Poglavlje 2} --- Regularni izrazi i njihova algebra.
  \item \textbf{Poglavlje 3} --- Konačni automati (KA, NKA),
        $\eps$-prijelazi, partitivna konstrukcija.
  \item \textbf{Poglavlje 4} --- Ekvivalentnost klasa regularnih jezika
        (RI, KA, NKA, DLG, LLG); Thompsonova konstrukcija, Ardenovo pravilo.
  \item \textbf{Poglavlje 5} --- Bezkontekstni jezici, potisni automati,
        Chomskyjeva normalna forma, CYK algoritam.
  \item \textbf{Poglavlje 6} --- Leme o pumpanju za regularne i
        bezkontekstne jezike.
  \item \textbf{Poglavlje 7} --- Kontekstni jezici, ograničeni automati,
        kontekstne i monotone gramatike.
  \item \textbf{Poglavlje 8} --- Turingovi strojevi, odlučivost, svođenje
        i dijagonalizacija.
\end{itemize}
\end{kljucno}

% ================================================================
\section*{Kako čitati skriptu}
% ================================================================

Tekst je organiziran tako da se može čitati linearno --- svaki pojam
uvodi se prije nego što se koristi. Ako već imate nešto pripreme,
slobodno preskačite ono što vam je poznato; ali ako pojma nemate odakle
krenuti, \textbf{počnite od početka}, najbolje s olovkom u ruci.

Vizualno je sve organizirano u obojene okvire pa se brzo prepoznaje
što se gleda:

\medskip
\noindent
\begin{tabular}{ll}
  \textcolor{defborder}{$\blacksquare$}~\textbf{Plavo (definicije)}
    & precizan, formalan iskaz pojma \\
  \textcolor{exborder}{$\blacksquare$}~\textbf{Zeleno (primjeri)}
    & konkretne ilustracije, ``vidi i probaj'' \\
  \textcolor{keyborder}{$\blacksquare$}~\textbf{Žuto (ključne ideje)}
    & ono što morate ponijeti sa sobom \\
  \textcolor{taskborder}{$\blacksquare$}~\textbf{Ljubičasto (zadaci)}
    & vježbe za samostalan rad \\
\end{tabular}

\medskip
\noindent
Na kraju svakog poglavlja stoji kratak \textbf{Sažetak} s tablicom
najvažnijih oznaka --- pogodno za brzo ponavljanje neposredno pred
ispit.

% ================================================================
\section*{O stilu i izvorima}
% ================================================================

Naglasak je na \emph{intuiciji} i \emph{povezanosti} pojmova, a ne na
potpunoj formalnoj striktnosti svakog dokaza. Tamo gdje su dokazi
tehnički zahtjevni, dane su \emph{skice} ili upute na izvore. Cilj
nije matematička virtuoznost --- cilj je da, kad vam netko pokaže
neki jezik, instinktivno znate razmišljati: \emph{regularan?
bezkontekstni? mogu li ga uopće prepoznati strojem?}

Glavni izvori bili su predavanja i službene bilješke kolegija, te
Sipserov \emph{Introduction to the Theory of Computation}.
Eventualne pogreške isključivo su moje --- ako ih primijetite,
slobodno javite na
\href{mailto:petar.prenc@gmail.com}{petar.prenc@gmail.com}.

\vspace{1.5em}
\begin{flushright}
\itshape Sretno!\\[0.3em]
Petar Prenc\\
Pula, \today
\end{flushright}

\cleardoublepage

% =========================================================
%  Sadrzaj
% =========================================================
\tableofcontents
\cleardoublepage

% =========================================================
%  Poglavlja
% =========================================================
% Chapter 1 — includable by main.tex
% Title converted to a chapter heading
\chapter{Jezici i osnovne operacije}
\label{chap:1}

\begin{quote}
\itshape
Sve što ćemo raditi u ovom kolegiju vrti se oko jednog jednostavnog pitanja:
pripada li neka riječ nekom jeziku? Ovo poglavlje postavlja temelje ---
definira što su uopće ,,riječi'' i ,,jezici'', i uvodi osnovne operacije s
kojima ćemo ih graditi i kombinirati.
\end{quote}


% ================================================================
\section{Tri pogleda na isti objekt}
% ================================================================

Zamislite funkciju $f : \mathbb{R} \to \mathbb{R}$. Možemo je opisati na
barem tri načina: kao skup uređenih parova $\{(x, f(x)) \mid x \in \mathbb{R}\}$,
kao formulu poput $f(x) = x^2 + 1$, ili kao program koji za svaki ulaz $x$
vraća odgovarajući izlaz. Sva tri opisa govore o istoj stvari, ali svaki naglašava
nešto drugačije i koristan je u različitim situacijama.

Formalni jezici bave se istim trikom, ali za jezike umjesto funkcija.
\textbf{Jezik} je skup --- statični popis riječi koje mu pripadaju.
\textbf{Gramatika} je dinamično pravilo koje taj skup \emph{generira},
kreće od početnog simbola i primjenjuje pravila prepisivanja.
\textbf{Automat} je idealizirani stroj koji \emph{prepoznaje} jezik čitajući
ulaznu riječ znak po znak i na kraju odlučuje: ,,da'' ili ,,ne''.

Ono što čini ovaj predmet zanimljivim jest to da između ova tri pogleda
postoje duboke i neočite veze. Na primjer, nije odmah jasno zašto bi za svaku
gramatiku određenog tipa \emph{nužno} postojao automat koji prepoznaje isti
jezik --- ali to je upravo ono što ćemo dokazivati.

Sve klase koje ćemo proučavati organizirane su u \textbf{Chomskyjevu hijerarhiju},
kojoj se vraćamo na kraju poglavlja.

% ================================================================
\section{Abecede i riječi}
% ================================================================

Prije nego što možemo govoriti o jezicima, moramo dogovoriti s kojim simbolima
radimo.

\begin{definicija}[Abeceda]
\textbf{Abeceda} je konačni, neprazni skup $\Sig$. Elemente abecede zovemo
\textbf{znakovima} i često označavamo malim grčkim slovima: $\alpha, \beta,
\gamma, \ldots$
\end{definicija}

Ništa posebno --- abeceda je samo dogovor o tome koji su nam simboli na
raspolaganju. Ograničenje da mora biti \emph{konačna} je ključno: svi naši
strojevi i gramatike imat će konačni opis, pa moramo krenuti od konačne
abecede.

\begin{primjer}[Nekoliko uobičajenih abeceda]
U praksi ćemo najčešće raditi s:
\begin{itemize}
  \item $\Sig = \{0, 1\}$ --- binarna abeceda, dovoljna za kodiranje čega god želimo
  \item $\Sig = \{a, b\}$ ili $\Sig = \{a, b, c\}$ --- mali primjeri za ilustracije
  \item $\Sig = \{0, 1, \ldots, 9\}$ --- decimalne znamenke
\end{itemize}
Abeceda ne mora biti skup slova ili znamenki --- može biti skup tokena, boja,
instrukcija ili bilo čega drugog. Bitno je samo da je konačna.
\end{primjer}

\begin{definicija}[Riječ i njezina duljina]
\textbf{Riječ} nad abecedom $\Sig$ je konačni niz znakova iz $\Sig$:
\[
  w = \alpha_1 \alpha_2 \cdots \alpha_n, \qquad \alpha_i \in \Sig.
\]
\textbf{Duljina} riječi $w$, pisana $|w|$, jest broj znakova u njoj.
Pišemo također $|w|_\alpha$ za broj pojavljivanja znaka $\alpha$ u $w$.
\end{definicija}

Posebno mjesto zauzima \textbf{prazna riječ} $\eps$ --- jedina riječ duljine 0.
To nije znak, nego prazan niz znakova, i vrijedi $|\eps| = 0$.

\begin{napomena}
Prazna riječ $\eps$ i prazan \emph{jezik} $\emptyset$ nisu isto.
Prazna riječ \emph{postoji} kao objekt, samo ne sadrži nijedan znak.
Prazan jezik je skup koji ne sadrži nijednu riječ --- pa onda ni $\eps$.
Ovu razliku je lako zamijeniti na ispitu, pazite.
\end{napomena}

Uobičajeno je poistovjećivati znak $\alpha \in \Sig$ s jednakovrijednom
jednoznakovnom riječi. Dakle, $a$ kao znak i $a$ kao jednoelementna
,,rečenica'' su isti objekt --- samo gledamo iz drugog kuta.

% ================================================================
\section{Konkatenacija}
% ================================================================

Osnovna operacija kojom iz kratkih riječi gradimo dulje jest
\textbf{konkatenacija} --- spajanje dvije riječi jedne za drugom.

\begin{definicija}[Konkatenacija]
Za riječi $u = \alpha_1 \cdots \alpha_n$ i $v = \beta_1 \cdots \beta_m$
definiramo
\[
  uv \;:=\; \alpha_1 \cdots \alpha_n \, \beta_1 \cdots \beta_m.
\]
Vrijedi $|uv| = |u| + |v|$ --- duljine se zbrajaju.
\end{definicija}

Konkatenacija je asocijativna: $(uv)w = u(vw)$, pa zagrade su suvišne
i pišemo jednostavno $uvw$. Prazna riječ $\eps$ je neutralni element:
$\eps w = w \eps = w$ --- dodavanje ničega ništa ne mijenja.

Budući da konkatenacija ,,izgleda kao množenje'', koristimo uobičajenu
multiplikativnu notaciju za potencije:
\[
  u^0 := \eps, \qquad u^1 := u, \qquad u^2 := uu, \qquad u^3 := uuu, \quad\ldots
\]

\begin{primjer}
Neka je $u = ab$ i $v = ba$. Tada $uv = abba$, $vu = baab$, $u^3 = ababab$.
Konkatenacija nije komutativna --- $uv \neq vu$ u općem slučaju.
\end{primjer}

% ================================================================
\section{Jezici}
% ================================================================

Sad kad znamo što su riječi, možemo definirati jezik.

\begin{definicija}[$\Sig^*$ i jezik]
Skup \emph{svih} mogućih riječi nad abecedom $\Sig$ (uključujući $\eps$)
označavamo s $\Sig^*$:
\[
  \Sig^* \;:=\; \{\eps\} \cup \Sig \cup \Sig^2 \cup \Sig^3 \cup \cdots
\]
\textbf{Jezik} nad $\Sig$ je bilo koji podskup $L \subseteq \Sig^*$.
\end{definicija}

Definicija je namjerno vrlo široka --- jezik je doslovno \emph{bilo koji}
skup riječi, konačni ili beskonačni. Jedino ograničenje je da su sve riječi
nad istom abecedom. Nema uvjeta da jezik mora imati neku posebnu strukturu;
to tek postaje zanimljivo kad počnemo klasificirati jezike po tome
\emph{kakvu} strukturu imaju.

\begin{primjer}
Nad $\Sig = \{a, b\}$, evo nekoliko jezika raznih vrsta:
\begin{itemize}
  \item $L_1 = \{a, ab, aba\}$ --- konačni jezik s tri elementa
  \item $L_2 = \{w \in \Sig^* \mid |w|_a = |w|_b\}$ --- sve riječi s jednakim
        brojem $a$-ova i $b$-ova; beskonačni
  \item $L_3 = \{a^n b^n \mid n \geq 1\} = \{ab, aabb, aaabbb, \ldots\}$ ---
        klasični primjer koji ćemo često susretati
  \item $\emptyset$ --- prazan jezik; nema nijedne riječi
  \item $\{\eps\}$ --- jezik s jednom riječi, praznom
  \item $\Sig^*$ --- svi mogući nizovi znakova, bez ikakva ograničenja
\end{itemize}
\end{primjer}

Jezike možemo kombinirati skupovnim operacijama ($\cup$, $\cap$,
$\setminus$, komplement spram $\Sig^*$), ali i operacijom konkatenacije,
proširenom s riječi na skupove:

\begin{definicija}[Konkatenacija jezika]
Za jezike $L, M \subseteq \Sig^*$ definiramo
\[
  LM \;:=\; \{uv \mid u \in L,\; v \in M\}.
\]
Dakle, $LM$ se dobiva lijepljenjem svake riječi iz $L$ sa svakom riječi iz $M$.
\end{definicija}

\begin{primjer}
Neka je $L = \{a, ab\}$ i $M = \{b, c\}$. Sve kombinacije su:
\[
  LM = \{a \cdot b,\; a \cdot c,\; ab \cdot b,\; ab \cdot c\}
     = \{ab,\; ac,\; abb,\; abc\}.
\]
A što je $L^2 = LL$? Svaka riječ iz $L$ zalijepljena sa svakom:
\[
  L^2 = \{aa,\; aab,\; aba,\; abab\}.
\]
Primijetite da $L^2$ ima točno $|L|^2 = 4$ elementa --- u ovom slučaju
nema ,,sudara'' (različitih kombinacija koje daju istu riječ). Nije
uvijek tako.
\end{primjer}

% ================================================================
\section{Kleeneve operacije}
% ================================================================

Najvažnija operacija u cijelom kolegiju je Kleeneva zvijezda.
Susresti ćete je svugdje --- u regularnim izrazima, gramatikama, posvuda.

\begin{definicija}[Kleenejev plus i zvijezda]
Za jezik $L \subseteq \Sig^*$ definiramo:
\begin{align*}
  L^+ &\;:=\; \bigcup_{k \geq 1} L^k
        \;=\; L \cup L^2 \cup L^3 \cup \cdots \\[6pt]
  L^* &\;:=\; \bigcup_{k \geq 0} L^k
        \;=\; \{\eps\} \cup L \cup L^2 \cup L^3 \cup \cdots
\end{align*}
Razlika je samo u tome što $L^*$ uključuje $L^0 = \{\eps\}$, a $L^+$ ne.
Vrijedi $L^* = \{\eps\} \cup L^+$, tj.\ zvijezda je plus uz dodatak
prazne riječi.
\end{definicija}

Kako čitati ove definicije? $L^*$ je skup svih \emph{nizanja} nula ili više
riječi iz $L$ jedne za drugom. Kad kažemo da neka riječ $w$ pripada $L^*$,
mislimo da se $w$ može ,,izrezati'' na komadiće pri čemu svaki komadić
pripada $L$.

\begin{kljucno}
$w \in L^*$ \textbf{ako i samo ako} postoji rastav
$w = w_1 w_2 \cdots w_k$ (za neko $k \geq 0$) takav da $w_i \in L$
za svaki $i$. Slučaj $k = 0$ daje $w = \eps$ --- uvijek prihvatljivo.
\end{kljucno}

\begin{primjer}
Uzmimo $L = \{0, 11\}$ nad $\Sig = \{0, 1\}$. Tada:
\[
  L^* = \{\eps,\; 0,\; 11,\; 00,\; 011,\; 110,\; 1111,\; 000,\; \ldots\}
\]
Je li $110011 \in L^*$? Da --- rastav je
$\underbrace{11}_{}\,\underbrace{0}_{}\,\underbrace{0}_{}\,\underbrace{11}_{}$,
svaki komadić je u $L$.

Je li $1 \in L^*$? Ne. Jedina jedinka u $L$ je $0$; jedini dvoznakac je $11$;
nema načina sastaviti samu $1$ od komadića iz $\{0, 11\}$.
\end{primjer}

\begin{primjer}
Posebno važan slučaj: ako je $\Sig$ abeceda (shvaćena kao jezik
jednoznakovnih riječi), tada je $\Sig^*$ upravo \emph{skup svih mogućih
riječi} nad $\Sig$ --- što je i definicija iz prethodnog odjeljka.
Zvijezda u oznaci $\Sig^*$ nije slučajnost.
\end{primjer}

Postoji korisna analogija s aritmetikom. Konkatenacija je množenje, prazna
riječ $\eps$ je jedinica, a Kleeneva zvijezda nalikuje na ,,beskonačnu
geometrijsku sumu'' $1 + L + L^2 + \cdots$, samo u algebri jezika.
Ta analogija nije samo poetska --- na njoj se temelji teorija regularnih
izraza koju ćemo graditi u sljedećem poglavlju.

\begin{zadatak}
\begin{enumerate}[label=(\alph*), itemsep=4pt]
  \item Izrazite $L^*$ pomoću $L^+$ i obrnuto. Koji uvjet mora biti ispunjen
        da vrijedi $L^+ = L^* \setminus \{\eps\}$?
  \item Dokažite ili opovrgnite: $(L \cup M)^* = (L^* M^*)^*$.
  \item Pokažite da je $\emptyset^* = \{\eps\}$.
        Intuitivno, zašto to ima smisla?
  \item Za $L = \{a\}$, opišite $L^*$ riječima. Koji je to dobro poznati skup?
\end{enumerate}
\end{zadatak}

% ================================================================
\section{Pozitivni jezici}
% ================================================================

Jedno naizgled sitno, ali tehnički korisno svojstvo jest pozitivnost.

\begin{definicija}[Pozitivan jezik]
Kažemo da je jezik $L$ \textbf{pozitivan} ako $\eps \notin L$.
\end{definicija}

Zašto nam je to uopće važno? U dokazima koji idu indukcijom po duljini
riječi, pozitivnost osigurava da konkatenacijom ne možemo ostati na istoj
duljini. Precizno: ako je $L$ pozitivan i $w \in LM$, dio koji pripada $M$
strogo je kraći od $w$ --- pa indukcija ,,napreduje''.

Vrijedi nekoliko jednostavnih pravila:
\begin{itemize}[itemsep=3pt]
  \item $\emptyset$ je pozitivan (trivijalno --- nema elemenata, pa ni $\eps$).
  \item Svaki jednoznakovni jezik $\{\alpha\}$ je pozitivan.
  \item $L^*$ \emph{nikad nije} pozitivan --- uvijek vrijedi $\eps \in L^*$.
  \item $LM$ je pozitivan ako i samo ako je barem jedan od $L$, $M$ pozitivan.
  \item $L \cup M$ je pozitivan ako i samo ako su \emph{oba} $L$ i $M$ pozitivna.
\end{itemize}

Zadnja dva pravila su dobar test razumijevanja --- probajte ih dokazati sami
iz definicija.

% ================================================================
\section{Chomskyjeva hijerarhija}
% ================================================================

Na kraju poglavlja smjestimo sve u širi okvir. Chomskyjeva hijerarhija
klasifikacija je jezika u klase prema tome koliko su ,,teški za prepoznati''.

Neformalna ideja: što je jezik strukturiraniji i zahtjevniji, to je moćniji
stroj (ili gramatika) potreban za njegovo opisivanje. Najjednostavniji jezici
--- \textbf{regularni} --- mogu se prepoznati \emph{bez ikakva pamćenja} osim
trenutnog stanja. Za malo zahtjevnije --- \textbf{bezkontekstne} --- treba nam
stog (neograničena memorija, ali pristup samo s vrha). Zatim dolaze
\textbf{kontekstni} jezici koji trebaju traku jednake duljine kao ulaz.
Na samom vrhu su \textbf{rekurzivno prebrojivi} jezici, za koje treba
puni Turingov stroj.

Svakoj klasi odgovara vrsta gramatike i vrsta automata:

\begin{center}
\renewcommand{\arraystretch}{1.5}
\begin{tabular}{llll}
\hline
\textbf{Klasa} & \textbf{Oznaka} & \textbf{Gramatika} & \textbf{Automat} \\
\hline
Regularni        & Reg & desno/lijevolinearna & konačni (KA, NKA) \\
Bezkontekstni    & BK  & bezkontekstna (BKG)  & potisni (PA) \\
Kontekstni       & Kon & kontekstna (KG)      & ograničeni (OA) \\
Rekurzivni       & Rek & ---                  & Turingov odlučitelj \\
Rek.\ prebrojivi & RE  & opća (OG)            & Turingov prepoznavač \\
\hline
\end{tabular}
\end{center}

Među klasama vrijede \emph{prave inkluzije} --- svaka klasa je pravi
podskup sljedeće:
\[
  \mathrm{Reg} \;\subsetneq\; \mathrm{BK} \;\subsetneq\; \mathrm{Kon}
  \;\subsetneq\; \mathrm{Rek} \;\subsetneq\; \mathrm{RE}.
\]
,,Prava'' inkluzija znači da za svaki par susjednih klasa postoji jezik koji
pripada gornjoj, ali ne i donjoj. Klasični primjeri koji ,,razdvajaju''
klase su:
\begin{align*}
  L_{=}      &= \{a^n b^n \mid n \geq 1\} \;\in\; \mathrm{BK} \setminus \mathrm{Reg}, \\
  L_{\equiv} &= \{a^n b^n c^n \mid n \geq 1\} \;\in\; \mathrm{Kon} \setminus \mathrm{BK}.
\end{align*}

Zašto, primjerice, $L_=$ nije regularan? Intuitivno: kad konačni automat
čita $a^n b^n$, u trenutku kad dođe do prvih $b$-ova mora nekako ,,znati''
koliko je $a$-ova pročitao --- a ima samo konačno mnogo stanja, pa ne može
pamtiti proizvoljan broj. Formalni dokaz dat ćemo u poglavlju o lemi o
pumpanju.

\begin{kljucno}[Što nas čeka]
Kolegij prolazi kroz hijerarhiju odozdo prema gore: regularni jezici
(poglavlja 2--4), bezkontekstni (poglavlje 5), kontekstni (poglavlje 7),
i na kraju Turingovi strojevi i odlučivost (poglavlje 8).
Svaki put kad pređemo u višu klasu, pronaći ćemo jezik koji ne stane
u onu ispod --- i naučiti \emph{zašto}.
\end{kljucno}

% ================================================================
\section*{Sažetak}
% ================================================================

Ovim smo postavili temelje. U jednoj rečenici: jezik je skup riječi nad
nekom abecedom, a ovo poglavlje definira kako se s takvim skupovima radi.

\medskip
\noindent
\begin{tabular}{ll}
  $\Sig$ & abeceda --- konačni, neprazni skup znakova \\
  $w, \eps$ & riječ, prazna riječ \\
  $|w|$ & duljina riječi \\
  $\Sig^*$ & skup svih riječi nad $\Sig$ \\
  $L \subseteq \Sig^*$ & jezik \\
  $LM$ & konkatenacija jezika \\
  $L^*$, $L^+$ & Kleeneva zvijezda i plus \\
\end{tabular}

\medskip
\noindent U sljedećem poglavlju uvodimo \textbf{regularne izraze} ---
kompaktan formalizam za opisivanje regularnih jezika koji koristi upravo
operacije $\cup$, konkatenaciju i ${}^*$ koje smo ovdje upoznali.

% Chapter 2 — includable by main.tex
\chapter{Regularni izrazi}
\label{chap:2}

\begin{quote}
\itshape
U prvom poglavlju naučili smo što su jezici i kako ih kombinirati.
Sad postavljamo pitanje: kako ih \emph{opisati}? Regularni izrazi daju
nam kompaktan, elegantan formalizam koji bi trebao izgledati poznato ---
vjerojatno ste već koristili \texttt{grep} ili \texttt{re} modul u Pythonu.
Ovo poglavlje tu intuiciju pretvara u preciznu matematiku.
\end{quote}


% ================================================================
\section{Motivacija: kako opisati beskonačni jezik?}
% ================================================================

Zamislite da vam netko kaže: ,,napiši mi sve valjane e-mail adrese''.
Odmah je jasno da taj skup ne možete nabrojati --- beskonačni je i nema
gornje granice duljine. No ipak intuitivno znate što je valjana adresa:
neka slova, pa \texttt{@}, pa domena. Regularni izrazi su upravo alat koji
takav opis pretvara u preciznu matematičku definiciju.

U Pythonu ste možda pisali nešto poput \verb|re.match(r'[a-z]+@[a-z]+\.[a-z]+', s)|.
Formalni regularni izrazi koje ćemo definirati su strožija i čišća verzija
iste ideje --- bez šarenih proširenja koja praktični jezici dodaju, ali s
jasnom matematičkom semantikom.

Ideja je jednostavna: regularni izrazi grade se od malih komadića
(\emph{inicijalnih} ili \emph{atomarnih} jezika) kombiniranjem s tri
operacije --- unije, konkatenacije i Kleeneve zvijezde. Ništa više.
Pokazat će se da je to i više nego dovoljno za cijelu klasu regularnih
jezika.

% ================================================================
\section{Definicija regularnih izraza}
% ================================================================

\begin{definicija}[RI-regularni jezik]
Neka je $\Sig$ abeceda. \textbf{Skup RI-regularnih jezika} nad $\Sig$
je \emph{najmanji} skup jezika koji:
\begin{itemize}
  \item sadrži $\emptyset$ i $\{\eps\}$ kao elemente,
  \item sadrži $\{\alpha\}$ za svaki znak $\alpha \in \Sig$,
  \item je zatvoren na konkatenaciju, konačnu uniju i Kleeneve operacije.
\end{itemize}
Jezik koji pripada tom skupu zovemo \textbf{RI-regularnim}.
\textbf{Regularni izraz} je izraz koji pokazuje kako je taj jezik izgrađen
od navedenih operacija.
\end{definicija}

Dakle, regularni izraz nije sam po sebi jezik --- on je \emph{recept} za
gradnju jezika. Kao što izraz $2 \cdot 3 + 1$ nije broj već opis kako
doći do broja $7$, regularni izraz $ab^*$ nije skup nego opis kako
doći do skupa $\{a, ab, abb, abbb, \ldots\}$.

\subsection{Inicijalni jezici}

Gradivni blokovi su sljedeća tri tipa:

\begin{center}
\renewcommand{\arraystretch}{1.5}
\begin{tabular}{cll}
\hline
\textbf{Izraz} & \textbf{Jezik} & \textbf{Opis} \\
\hline
$\emptyset$ & $\emptyset$ & prazan jezik, nema nijedne riječi \\
$\eps$      & $\{\eps\}$  & samo prazna riječ \\
$\alpha$    & $\{\alpha\}$ & samo jednoznakovna riječ od $\alpha$ \\
\hline
\end{tabular}
\end{center}

\begin{napomena}
Pažnja na oznake: $\emptyset$ kao regularni izraz opisuje \emph{prazan jezik},
a $\eps$ kao regularni izraz opisuje jezik $\{\eps\}$ koji sadrži \emph{jednu}
riječ. Ovo je ista razlika kao u prethodnom poglavlju --- samo sad gledamo
izraze, ne skupove.
\end{napomena}

\subsection{Operacije}

Iz inicijalnih jezika gradimo veće regularnim operacijama:

\begin{center}
\renewcommand{\arraystretch}{1.5}
\begin{tabular}{cll}
\hline
\textbf{Izraz} & \textbf{Jezik} & \textbf{Čitamo} \\
\hline
$r \cup s$  & $L(r) \cup L(s)$   & unija: ili $r$ ili $s$ \\
$r \cdot s$ & $L(r) \cdot L(s)$  & konkatenacija: prvo $r$, pa $s$ \\
$r^*$       & $L(r)^*$           & zvijezda: nula ili više puta $r$ \\
$r^+$       & $L(r)^+$           & plus: jedanput ili više puta $r$ \\
\hline
\end{tabular}
\end{center}

Točka za konkatenaciju često se izostavlja: $rs$ znači isto što i $r \cdot s$.
Prioritet operacija je kao u aritmetici: ${}^*$ i ${}^+$ vežu najjače, pa
konkatenacija, pa unija. Dakle $ab^*c \cup d$ čitamo kao
$(a \cdot (b^*) \cdot c) \cup d$.

% ================================================================
\section{Primjeri --- od jednostavnog prema složenom}
% ================================================================

Sad kad imamo definiciju, vježbamo je na konkretnim primjerima. Ključna
vještina je dvojaka: čitati regularni izraz i razumjeti koji jezik opisuje,
i zadani jezik opisati regularnim izrazom.

\begin{primjer}[Jednostavni izrazi]
Nad $\Sig = \{a, b\}$:
\begin{itemize}
  \item $a \cup b$ opisuje jezik $\{a, b\}$ --- jedna od ta dva slova.
  \item $ab$ opisuje $\{ab\}$ --- točno ta jedna dvoznakovna riječ.
  \item $a^*$ opisuje $\{\eps, a, aa, aaa, \ldots\}$ --- nula ili više $a$-ova.
  \item $a^+$ opisuje $\{a, aa, aaa, \ldots\}$ --- jedan ili više $a$-ova.
  \item $a^* b$ opisuje $\{b, ab, aab, aaab, \ldots\}$ --- nula ili više $a$-ova, pa jedan $b$.
\end{itemize}
\end{primjer}

\begin{primjer}[Sve riječi nad $\{a, b\}$]
Izraz $(a \cup b)^*$ opisuje \emph{sve} moguće riječi nad $\{a, b\}$,
dakle $\Sig^*$. Kleeneva zvijezda od $\{a\} \cup \{b\} = \{a,b\}$
daje sve načine nizanja $a$-ova i $b$-ova, bez ikakva ograničenja.

To je korisna crtica: kad god trebate ,,bilo što od tih znakova'',
pišete $(a \cup b)^*$, ili općenito $\Sig^*$.
\end{primjer}

\begin{primjer}[Binarni brojevi]
Razmotrimo izraz $1(0 \cup 1)^* \cup 0$ nad $\Sig = \{0, 1\}$.
\begin{itemize}
  \item $1(0 \cup 1)^*$: počinje s $1$, a zatim slijedi bilo što ---
        to su svi binarni brojevi $\geq 1$.
  \item $0$: baš sama nula.
  \item Unija: skupa dobivamo \emph{sve valjane binarne zapise prirodnih
        brojeva} bez vodećih nula (jedino $0$ počinje s nulom).
\end{itemize}
Ovo je tipičan uzorak: posebnim slučajem ($0$) riješimo rub,
a glavnom granom ($1(0\cup1)^*$) pokrijemo opći slučaj.
\end{primjer}

\begin{primjer}[Riječi koje počinju i završavaju istim slovom]
Napišimo regularni izraz za sve riječi nad $\{a, b, c\}$ koje
počinju i završavaju istim slovom. Koristimo kraticu $\Sig := a \cup b \cup c$.

Fiksiramo koje je to slovo i ujedinimo tri slučaja:
\begin{itemize}
  \item Počinje i završava s $a$: treba $a \Sig^* a$, ali i samu riječ $a$.
        Dakle: $a(\Sig^* a \cup \eps)$.
  \item Analogno za $b$ i $c$.
\end{itemize}
Ukupno:
\[
  a(\Sig^* a \cup \eps) \;\cup\; b(\Sig^* b \cup \eps) \;\cup\; c(\Sig^* c \cup \eps).
\]
\end{primjer}

\begin{primjer}[Neparan broj $a$-ova]
Jezik svih riječi nad $\{a,b\}$ s neparnim brojem $a$-ova.

Intuicija: između $a$-ova može biti koliko god $b$-ova. Pokušajmo
$b^* (ab^*)^+$ --- ali to dopušta \emph{bilo koji} broj $a$-ova, parni ili neparni.

Rješenje: svaki par $a$-ova se ,,poništi'', a jedan ostaje:
\[
  b^* \bigl(ab^* ab^*\bigr)^* ab^*.
\]
Čitanje: nula ili više parova $a \cdots a$ (okruženih $b$-ovima),
pa točno jedan $a$, pa $b$-ovi. Ukupan broj $a$-ova je $2k + 1$ za neko
$k \geq 0$ --- uvijek neparan.
\end{primjer}

\begin{kljucno}
Izgradnja regularnog izraza za zadani jezik obično ide ovako:
\begin{enumerate}
  \item Identificirajte strukturu riječi --- što dolazi na kojoj poziciji.
  \item Podijelite na slučajeve ako treba (npr.\ različiti početni znakovi).
  \item Koristite $(a \cup b \cup \ldots)^*$ za ,,bilo što od tih znakova''.
  \item Provjerite rubne slučajeve: prazna riječ, jednoznakovne riječi.
\end{enumerate}
\end{kljucno}

% ================================================================
\section{Algebarske jednakosti}
% ================================================================

Dva regularna izraza $r$ i $s$ su \textbf{ekvivalentni} ako opisuju isti
jezik: $L(r) = L(s)$. Pišemo $r \equiv s$.

Postoji niz korisnih algebarskih identiteta:

\begin{center}
\renewcommand{\arraystretch}{1.6}
\begin{tabular}{ll}
\hline
\textbf{Identitet} & \textbf{Što kaže} \\
\hline
$r \cup \emptyset \equiv r$ & unija s praznim jezikom ne mijenja ništa \\
$r \cdot \eps \equiv r$     & konkatenacija s $\eps$ ne mijenja ništa \\
$r \cdot \emptyset \equiv \emptyset$ & konkatenacija s praznim daje prazno \\
$r \cup r \equiv r$         & unija sa samim sobom je idempotentna \\
$r \cup s \equiv s \cup r$  & unija je komutativna \\
$(rs)t \equiv r(st)$        & konkatenacija je asocijativna \\
$r(s \cup t) \equiv rs \cup rt$ & distributivnost \\
$(r^*)^* \equiv r^*$        & zvijezda od zvijezde je zvijezda \\
$\eps^* \equiv \eps$        & zvijezda od $\eps$ je $\eps$ \\
$\emptyset^* \equiv \eps$   & zvijezda od praznog skupa je $\{\eps\}$ \\
\hline
\end{tabular}
\end{center}

\begin{primjer}
Provjerimo $\emptyset^* \equiv \eps$, tj.\ $L(\emptyset^*) = \{\eps\}$.
Po definiciji, $\emptyset^* = \bigcup_{k \geq 0} \emptyset^k$.
Vrijedi $\emptyset^0 = \{\eps\}$ i $\emptyset^k = \emptyset$ za $k \geq 1$.
Dakle $\emptyset^* = \{\eps\}$. Intuitivno: jedini način da se ,,niže nula
ili više riječi iz praznog skupa'' je uzeti nula riječi, što daje $\eps$.
\end{primjer}

\begin{zadatak}
\begin{enumerate}[label=(\alph*), itemsep=4pt]
  \item Pokažite da je $r^+ \equiv rr^* \equiv r^*r$.
  \item Vrijedi li $r^* \equiv r^+ \cup \eps$ za svaki $r$?
        Što ako je $\eps \in L(r)$?
  \item Nad $\Sig = \{0,1\}$, napišite regularni izraz za:
    \begin{itemize}
      \item sve riječi koje završavaju na $01$,
      \item sve riječi koje sadrže podniz $11$,
      \item sve riječi parne duljine.
    \end{itemize}
  \item Nad $\Sig = \{a, b\}$, napišite regularni izraz za sve riječi
        u kojima se $a$ nikad ne pojavljuje dva puta zaredom.
\end{enumerate}
\end{zadatak}

% ================================================================
\section{Što regularni izrazi ne mogu?}
% ================================================================

Regularni izrazi su moćni, ali nisu svemoćni.

Jezik $L_= = \{a^n b^n \mid n \geq 1\}$ sadrži riječi $ab$, $aabb$,
$aaabbb$, \ldots\ Svaka riječ ima isti broj $a$-ova i $b$-ova.
Pokušajte napisati regularni izraz za taj jezik --- neće ići.

Zašto? Regularni izraz ne može ,,prebrojati'' i usporediti duljine dva
dijela riječi. Sve što može je pratiti konačno mnogo mogućih ,,stanja''
--- a mogućih vrijednosti $n$ je beskonačno mnogo. Formalni dokaz dolazi
s \textbf{lemom o pumpanju} (poglavlje~6), ali intuicija je jasna već sad.

\begin{kljucno}
Regularni izrazi opisuju jezike u kojima se struktura može pratiti bez
pamćenja neograničenih brojača. Čim trebate ,,zapamtiti'' koliko puta
ste nešto vidjeli i usporediti s nekim drugim brojem --- regularni izraz
nije dovoljan alat.
\end{kljucno}

% ================================================================
\section*{Sažetak}
% ================================================================

Regularni izrazi daju nam kompaktan način opisivanja regularnih jezika.
Grade se od tri atomarna slučaja ($\emptyset$, $\eps$, jednoznakovni jezici)
i tri operacije (unija, konkatenacija, Kleeneva zvijezda/plus).

\medskip
\noindent
\begin{tabular}{ll}
  $\emptyset$  & prazan jezik \\
  $\eps$       & jezik $\{\eps\}$ \\
  $\alpha$     & jezik $\{\alpha\}$ \\
  $r \cup s$   & unija \\
  $rs$         & konkatenacija \\
  $r^*$        & Kleeneva zvijezda (nula ili više puta) \\
  $r^+$        & Kleenejev plus (jedanput ili više) \\
\end{tabular}

\medskip
\noindent U sljedećem poglavlju uvodimo \textbf{konačne automate} ---
strojeve koji prepoznaju regularne jezike. Vidjet ćemo da su automati i
regularni izrazi \emph{ekvivalentni}: svaki regularni izraz daje automat,
i svaki automat daje regularni izraz.

% Chapter 3 — includable by main.tex
\chapter{Konačni automati}
\label{chap:3}

\begin{quote}
\itshape
Regularni izrazi opisuju što jezik sadrži --- ali kako provjeriti
pripada li konkretna riječ nekom jeziku? Za to trebamo stroj.
Konačni automat je najjednostavniji takav stroj: čita riječ znak po
znak i na kraju odlučuje da ili ne, bez ikakve memorije osim
trenutnog stanja. Ovaj jednostavni model ispostavlja se jednako
moćnim kao i regularni izrazi.
\end{quote}


% ================================================================
\section{Ideja: stroj koji čita riječ}
% ================================================================

Zamislite jednostavan sustav za provjeru PIN-a bankomata. On prima
niz znamenki i na kraju kaže: točno ili pogrešno. U svakom trenutku
pamti samo ``koliko smo znakova dosad prihvatili'' --- nema nikakve
druge memorije. Takav sustav je, u biti, konačni automat.

Formalno, konačni automat je idealizirani matematički stroj koji:
\begin{itemize}
  \item ima konačno mnogo \textbf{stanja} (kao moguće ``faze'' rada),
  \item kreće iz fiksnog \textbf{početnog stanja},
  \item za svaki pročitani znak prelazi u novo stanje prema
        \textbf{funkciji prijelaza},
    \item na kraju prihvaća riječ ako je završio u nekom od
      \textbf{završnih stanja}.
\end{itemize}

Kljucno ogranicenje: nema stoga, nema trake, nema ekstre memorije.
Jedino sto stroj ``zna'' jest u kojem je stanju --- i to je sve.

% ================================================================
\section{Deterministički konačni automat (KA)}
% ================================================================

\begin{definicija}[Konačni automat]
\textbf{Konačni automat} (KA) nad abecedom $\Sig$ je uređena petorka
\[
  M = (Q,\, \Sig,\, \delta,\, q_0,\, F),
\]
gdje je:
\begin{itemize}
  \item $Q$ --- konačni skup \textbf{stanja},
  \item $\Sig$ --- abeceda (konačni skup ulaznih znakova),
  \item $\delta : Q \times \Sig \to Q$ --- \textbf{funkcija prijelaza},
  \item $q_0 \in Q$ --- \textbf{početno stanje},
  \item $F \subseteq Q$ --- skup \textbf{završnih stanja}.
\end{itemize}
\end{definicija}

Funkcija prijelaza $\delta$ je srce automata: za svako stanje i svaki
ulazni znak precizno govori u koje stanje automat prelazi. Nema
nedefiniranih prijelaza, nema izbora --- za svaki par $(q, \alpha)$
postoji točno jedan nasljednik $\delta(q, \alpha)$.

\begin{definicija}[Prihvaćanje riječi]
KA $M = (Q, \Sig, \delta, q_0, F)$ \textbf{prihvaća} riječ
$w = \alpha_1 \alpha_2 \cdots \alpha_n \in \Sig^*$ ako postoji niz stanja
$r_0, r_1, \ldots, r_n \in Q$ takav da:
\begin{itemize}
  \item $r_0 = q_0$ (kreće iz početnog stanja),
  \item $\delta(r_{i-1}, \alpha_i) = r_i$ za sve $1 \leq i \leq n$ (prati prijelaze),
  \item $r_n \in F$ (završava u završnom stanju).
\end{itemize}
Jezik koji $M$ prepoznaje je
\[
  L(M) := \{w \in \Sig^* : M \text{ prihvaća } w\}.
\]
\end{definicija}

\begin{napomena}
Za deterministički KA, niz stanja $r_0, r_1, \ldots, r_n$ je
\emph{jedinstven} za zadane $M$ i $w$ --- nema izbora ni slučajnosti.
Automat uvijek "zna" što radi.
\end{napomena}

\begin{primjer}[KA koji prihvaća riječi s parnim brojem $a$-ova]
Nad $\Sig = \{a, b\}$, konstruirajmo KA za jezik svih riječi s
parnim brojem pojavljivanja slova $a$ (uključujući $0$).

Trebamo pratiti samo parnost broja $a$-ova --- dvije mogucnosti:
``parno'' i ``neparno''. To su nase dvije stanice:
\[
  Q = \{q_{\text{par}},\, q_{\text{nepar}}\},
  \quad q_0 = q_{\text{par}},
  \quad F = \{q_{\text{par}}\}.
\]
Funkcija prijelaza:
\begin{center}
\renewcommand{\arraystretch}{1.4}
\begin{tabular}{c|cc}
  & $a$ & $b$ \\ \hline
  $q_{\text{par}}$   & $q_{\text{nepar}}$ & $q_{\text{par}}$ \\
  $q_{\text{nepar}}$ & $q_{\text{par}}$   & $q_{\text{nepar}}$ \\
\end{tabular}
\end{center}
Riječ $aabba$: $q_{\text{par}} \xrightarrow{a} q_{\text{nepar}}
\xrightarrow{a} q_{\text{par}} \xrightarrow{b} q_{\text{par}}
\xrightarrow{b} q_{\text{par}} \xrightarrow{a} q_{\text{nepar}}$.
Završava u $q_{\text{nepar}} \notin F$ --- riječ nije prihvaćena (ima $3$ slova $a$). \checkmark
\end{primjer}

\begin{primjer}[KA koji prihvaća riječi koje završavaju na $ab$]
Nad $\Sig = \{a, b\}$, želimo prepoznati sve riječi koje završavaju
nizom $ab$. Automat mora pamtiti ``suffix koji smo dosad vidjeli'',
ali samo onoliko koliko je potrebno:
\[
  Q = \{q_0,\, q_a,\, q_{ab}\}, \quad F = \{q_{ab}\}.
\]
\begin{center}
\renewcommand{\arraystretch}{1.4}
\begin{tabular}{c|cc}
  & $a$ & $b$ \\ \hline
  $q_0$  & $q_a$  & $q_0$ \\
  $q_a$  & $q_a$  & $q_{ab}$ \\
  $q_{ab}$ & $q_a$ & $q_0$ \\
\end{tabular}
\end{center}
Stanje $q_a$ znači: ``zadnji pročitani znak je bio $a$''. Stanje $q_{ab}$:
``zadnja dva znaka su bila $ab$''. Iz $q_{ab}$, ako pročitamo $a$, idemo
u $q_a$ (novi potencijalni početak); ako pročitamo $b$, gubimo trag
i vraćamo se u $q_0$.
\end{primjer}

% ================================================================
\section{Nedeterministički konačni automat (NKA)}
% ================================================================

Deterministički KA je konceptualno jasan, ali često nespretan za
dizajnirati. Prirodniji pristup je često \emph{nedeterminističan}:
u svakom trenutku automat smije ``pogoditi'' koji će prijelaz napraviti,
a prihvaća riječ ako \emph{postoji barem jedna grana} koja vodi u
završno stanje.

\begin{definicija}[Nedeterministički KA]
\textbf{Nedeterministički konačni automat} (NKA) nad $\Sig$ je
uređena petorka $(Q, \Sig, \Delta, q_0, F)$, gdje je sve kao kod KA
osim sto je funkcija prijelaza zamijenjena \textbf{relacijom prijelaza}:
\[
  \Delta \subseteq Q \times \Sig_\eps \times Q,
\]
gdje je $\Sig_\eps := \Sig \cup \{\eps\}$. Prijelaz $(q, \alpha, p) \in \Delta$
znači: ``iz stanja $q$, čitajući znak $\alpha$ (ili, ako je $\alpha = \eps$,
bez čitanja), može se preći u stanje $p$''.

NKA prihvaća riječ $w$ ako postoji \emph{barem jedan} način čitanja
$w$ koji vodi iz $q_0$ u neko završno stanje.
\end{definicija}

\begin{kljucno}
Nedeterminizam je matematički trik: NKA ``pogađa'' točnu granu.
U stvarnosti takav stroj ne postoji --- ali matematički je koristan
jer je često mnogo lakši za dizajnirati od ekvivalentnog KA.
Ključni rezultat: svaki NKA ima ekvivalentni KA!
\end{kljucno}

\subsection{Epsilon-prijelazi}

Posebno su korisni \textbf{$\eps$-prijelazi}: automat prelazi iz jednog
stanja u drugo \emph{bez čitanja ijednog znaka}. Ovo se čini kao
varanje, ali zapravo samo kaže: ``automat smije besplatno proći kroz
neka stanja''.

\begin{definicija}[$\eps$-ljuska]
Za skup stanja $S \subseteq Q$, \textbf{$\eps$-ljuska} od $S$ je skup
svih stanja dostupnih iz $S$ iskljucivo $\eps$-prijelazima:
\[
  [S]_\eps := \{q' \in Q : \text{postoji staza } \eps\text{-prijelaza od nekog } q \in S \text{ do } q'\}.
\]
\end{definicija}

\begin{primjer}[$\eps$-prijelazi u praksi]
Kad želimo spojiti dva NKA $N_1$ i $N_2$ tako da prepoznaju
konkatenaciju $L(N_1) L(N_2)$, to radimo tako da iz svakog završnog
stanja od $N_1$ dodamo $\eps$-prijelaz prema početnom stanju od $N_2$.
Niti jedan znak se ne ``gubi'' --- samo automati postaju lijepo spojeni.
\end{primjer}

% ================================================================
\section{Ekvivalentnost KA i NKA}
% ================================================================

\begin{propozicija}[Sipser, teorem 1.39]
Jezik je KA-regularan ako i samo ako je NKA-regularan.
\end{propozicija}

Smjer $(\Rightarrow)$ je trivijalan: svaki KA je specijalan slucaj NKA
(bez nedeterminizma i $\eps$-prijelaza). Smjer $(\Leftarrow)$ je
zanimljiviji i koristi \textbf{partitivnu konstrukciju}.

\subsection{Partitivna konstrukcija}

Ideja: ako NKA u svakom trenutku može biti u više stanja odjednom,
zasto ne bismo cijeli skup mogućih stanja proglasili jednim stanjem
novog, determinističkog automata?

\begin{definicija}[Partitivna konstrukcija]
Neka je $N = (Q, \Sig, \Delta, q_0, F)$ NKA. Definiramo KA
\[
  P(N) = \bigl(\mathcal{P}(Q),\; \Sig,\; \delta,\; [\{q_0\}]_\eps,\;
  \{T \subseteq Q : T \cap F \neq \emptyset\}\bigr),
\]
gdje je
\[
  \delta(T, \alpha) := \bigl[\{q \in Q : (\exists p \in T)\; (p, \alpha, q) \in \Delta\}\bigr]_\eps.
\]
\end{definicija}

Intuicija: stanje $P(N)$-a je skup stanja u kojima $N$ \emph{može}
biti. Kad pročitamo znak $\alpha$, uzimamo sve $\Delta$-nasljednike
svakog stanja iz skupa, a zatim zatvorimo na $\eps$-prijelaze.

\begin{napomena}
Ako $N$ ima $n$ stanja, $P(N)$ ima $2^n$ stanja --- eksponencijalan
rast! U praksi se većinom koriste samo dostupna stanja (ona do kojih
se uopće može doći iz početnog), kojih je obično znatno manje.
\end{napomena}

\begin{primjer}[Partitivna konstrukcija]
Neka je $N$ NKA s $Q = \{1, 2, 3\}$, $q_0 = 1$, $F = \{3\}$ i
prijelazima:
\[
  \Delta = \{(1,a,1),\; (1,a,2),\; (2,b,3),\; (3,b,3)\}.
\]
Ovaj NKA prepoznaje rijeci oblika $a^n b^k$ za $n \geq 1$, $k \geq 1$
(nejasno, ali pogledajmo partitivnu konstrukciju):

Početno stanje $P(N)$-a je $\{1\}$. Pratimo:
\begin{itemize}
  \item $\delta(\{1\}, a) = \{1, 2\}$ (jer $(1,a,1)$ i $(1,a,2)$)
  \item $\delta(\{1,2\}, a) = \{1, 2\}$ (iz 1 kao gore, iz 2 nema)
  \item $\delta(\{1,2\}, b) = \{3\}$ (samo $(2,b,3)$)
  \item $\delta(\{3\}, b) = \{3\}$ i $\delta(\{3\}, a) = \emptyset$
\end{itemize}
Dostupna stanja: $\{1\}$, $\{1,2\}$, $\{3\}$, $\emptyset$ ---
samo $4$ stanja umjesto $2^3 = 8$.
\end{primjer}

% ================================================================
\section{Zatvorenost na skupovne i jezicne operacije}
% ================================================================

Jedna od ključnih osobina regularnih jezika jest da su zatvoreni na
mnoge prirodne operacije.

\subsection{Skupovne operacije}

\begin{propozicija}
Klasa KA-regularnih jezika zatvorena je na \textbf{uniju}, \textbf{presjek}
i \textbf{komplement}.
\end{propozicija}

Dokaz za uniju i presjek koristi \textbf{Kartezijevu konstrukciju}:
ako $M_1$ prepoznaje $L_1$ i $M_2$ prepoznaje $L_2$, konstruiramo
KA koji ``istovremeno'' simulira oba:
\[
  M_1 \times M_2 := (Q_1 \times Q_2,\; \Sig,\; \delta,\; (q_1, q_2),\; F),
\]
gdje $\delta((q,q'), \alpha) := (\delta_1(q,\alpha), \delta_2(q',\alpha))$.

Za \textbf{uniju} stavimo $F = Q_1 \times F_2 \cup F_1 \times Q_2$.
Za \textbf{presjek} stavimo $F = F_1 \times F_2$.
Za \textbf{komplement} samo zamijenimo završna s nezavršnima: $F' = Q \setminus F$.

\subsection{Jezicne operacije}

\begin{propozicija}
Klasa NKA-regularnih (= regularnih) jezika zatvorena je na
\textbf{konkatenaciju}, \textbf{Kleenevu zvijezdu} i \textbf{Kleenejev plus}.
\end{propozicija}

Ove konstrukcije su elegantnije s NKA:
\begin{itemize}
    \item \textbf{Unija $L \cup L'$}: dodamo novo početno stanje s
      $\eps$-prijelazima prema početnim stanjima oba automata.
    \item \textbf{Konkatenacija $LL'$}: iz svakog završnog stanja
      od $N$ dodamo $\eps$-prijelaz prema početnom stanju od $N'$.
      Stara završna stanja od $N$ više nisu završna.
    \item \textbf{Kleeneva zvijezda $L^*$}: dodamo novo početno završno
      stanje s $\eps$-prijelazima prema starim početnim stanjima,
      te iz svakog završnog stanja $\eps$-prijelaz natrag prema
      početnom stanju.
\end{itemize}

% ================================================================
\section{Kako dizajnirati KA: heuristike}
% ================================================================

Dizajn KA je vjestina koja se uci vjezbanjem. Nekoliko heuristika:

\begin{kljucno}
\textbf{Pitajte se: sto KA mora pamtiti?}

Stanja KA su različite ``situacije'' u kojima se automat može naći.
Pitanje je: koje informacije o dosad pročitanom prefiksu su relevantne
za odluku o prihvaćanju?

Primjeri:
\begin{itemize}
  \item ``Zadnji pročitani znak'' --- $|\Sig|$ stanja.
  \item ``Parnost broja $a$-ova'' --- 2 stanja.
  \item ``Ostatak duljine modulo $k$'' --- $k$ stanja.
  \item ``Prefiks duljine do $k$'' --- konačno mnogo stanja.
\end{itemize}
Ako trebate pamtiti nešto \emph{neograničeno} (npr.\ točni broj $a$-ova)
--- KA za to ne postoji.
\end{kljucno}

\begin{zadatak}
\begin{enumerate}[label=(\alph*), itemsep=4pt]
  \item Konstruirajte KA nad $\{a,b\}$ koji prepoznaje jezik svih rijeci
        cija duljina je djeljiva s $3$.
  \item Konstruirajte NKA koji prepoznaje jezik svih rijeci nad $\{0,1\}$
        koje završavaju s $010$.
  \item Primijenite partitivnu konstrukciju na NKA iz (b) da dobijete
        ekvivalentni KA.
  \item Konstruirajte KA koji prepoznaje komplement jezika iz (a).
\end{enumerate}
\end{zadatak}

% ================================================================
\section*{Sazetak}
% ================================================================

Konačni automat čita riječ znak po znak, ne može se vraćati unatrag,
i pamti samo stanje u kojemu je. Usprkos ovim jakim ograničenjima,
konačni automati prepoznaju točno regularnu klasu jezika.

\medskip
\noindent
\begin{tabular}{ll}
  $M = (Q, \Sig, \delta, q_0, F)$ & deterministički KA \\
  $N = (Q, \Sig, \Delta, q_0, F)$ & nedeterministički NKA \\
  $[S]_\eps$ & $\eps$-ljuska skupa stanja \\
  $P(N)$ & partitivna konstrukcija (NKA $\to$ KA) \\
  $M_1 \times M_2$ & Kartezijeva konstrukcija (unija/presjek) \\
\end{tabular}

\medskip
\noindent U sljedećem poglavlju zatvaramo krug: pokazujemo da su
KA, NKA, regularni izrazi i linearne gramatike sve \emph{isti} formalizam,
samo u različitom ruhu.


% Chapter 4 — includable by main.tex
\chapter{Ekvivalentnost klasa regularnih jezika}
\label{chap:4}

\begin{quote}
\itshape
U dosad uveli nekoliko načina opisivanja regularnih jezika, a ovo
poglavlje zatvara krug — pokazat ćemo ekvivalentnost različitih formalizama.
\end{quote}

\newpage

% ================================================================
\section{Krug koji treba zatvoriti}
% ================================================================

Dosad smo uveli nekoliko načina opisivanja regularnih jezika:
\begin{itemize}
  \item \textbf{RI} --- regularni izrazi (poglavlje 2)
  \item \textbf{KA} --- deterministički konačni automati (poglavlje 3)
  \item \textbf{NKA} --- nedeterministički konačni automati (poglavlje 3)
  \item \textbf{DLG} --- desnolinearne gramatike (ovo poglavlje)
  \item \textbf{LLG} --- lijevolinearne gramatike (ovo poglavlje)
\end{itemize}

Svaki od ovih formalizma definira neku klasu jezika. Ključni teorem
ovog poglavlja kaže da su sve te klase iste:
\[
  \mathrm{Reg_{KA}} = \mathrm{Reg_{NKA}} = \mathrm{Reg_{RI}}
  = \mathrm{Reg_{DLG}} = \mathrm{Reg_{LLG}}.
\]
Tu klasu jednostavno zovemo $\mathrm{Reg}$ --- \textbf{klasa regularnih
jezika}.

Dokaz ide ``u krug'': pokazujemo $\mathrm{RI} \to \mathrm{NKA} \to
\mathrm{KA} \to \mathrm{DLG} \to \mathrm{RI}$, a $\mathrm{LLG}$
uključujemo pomoću operacije reverza.

% ================================================================
\section{Desnolinearne gramatike (DLG)}
% ================================================================

\begin{definicija}[Gramatika]
\textbf{Gramatika} nad abecedom $\Sig$ je transformacijski sustav
$G = (V, \Sig, \to, S)$, gdje je:
\begin{itemize}
  \item $V$ --- konačni skup \textbf{varijabli} (disjunktan od $\Sig$),
  \item $Y := \Sig \cup V$ --- skup svih \textbf{simbola},
  \item $S \in V$ --- \textbf{početna varijabla},
  \item ${\to} \subseteq Y^+ \times Y^*$ --- konačna relacija
        \textbf{pravila} (na lijevoj strani barem jedna varijabla).
\end{itemize}
Pisemo $u \Rightarrow v$ ako postoji pravilo $s \to t$ i rijeci $a, b$
takve da je $u = asb$ i $v = atb$. Gramatika \textbf{generira} jezik
\[
  L(G) := \{w \in \Sig^* : S \Rightarrow^* w\}.
\]
\end{definicija}

Ovo je opća definicija gramatike. Za regularne jezike dovoljne su nam
gramatike s \emph{jako} ograničenim pravilima:

\begin{definicija}[Desnolinearna gramatika]
Gramatika $G$ je \textbf{desnolinearna} (DLG) ako su sva njena pravila
jednog od dva oblika:
\[
  A \to \alpha B \qquad \text{ili} \qquad A \to \eps,
\]
gdje su $A, B \in V$ varijable, a $\alpha \in \Sig$ znak.
\end{definicija}

Intuicija: svako pravilo ili ``pojede'' jedan znak i zamijeni varijablu
s novom varijablom (koja će pojesti sljedeći), ili varijabla nestane
(kraj izvoda). Varijabla je uvijek na \emph{desnom} kraju međurječita
--- otuda ime ``desnolinearna''.

\begin{primjer}
DLG s varijablama $\{S, A\}$ i pravilima
$S \to aA$, $A \to aA$, $A \to bA$, $A \to \eps$
generira jezik svih rijeci koje pocinju s $a$ i zatim mogu imati bilo
sto. Regularni izraz: $a(a \cup b)^*$.
\end{primjer}

% ================================================================
\section{Pretvorba KA $\to$ DLG}
% ================================================================

Svaki KA mozemo pretvoriti u ekvivalentnu DLG na prirodan nacin:
stanja postaju varijable, prijelazi postaju pravila.

\begin{propozicija}
Svaki KA-regularni jezik je DLG-regularan.
\end{propozicija}

\textit{Dokaz (konstrukcija).}
Neka je $M = (Q, \Sig, \delta, q_0, F)$ KA koji prepoznaje jezik $L$.
Konstruiramo DLG $G$ ovako:
\begin{itemize}
  \item Za svako stanje $q \in Q$ uvedemo varijablu $V_q$.
  \item Početna varijabla: $S := V_{q_0}$.
  \item Za svaki prijelaz $\delta(q, \alpha) = p$ dodamo pravilo
        $V_q \to \alpha V_p$.
  \item Za svako završno stanje $q \in F$ dodamo pravilo $V_q \to \eps$.
\end{itemize}
Moze se pokazati indukcijom da za svako stanje $q \in Q$ vrijedi
\[
  L(Q, \Sig, \delta, q, F) = L(V, \Sig, \to, V_q),
\]
tj.\ varijabla $V_q$ generira tocno rijeci koje KA prihvaca krenuvsi
iz stanja $q$. Za $q := q_0$ dobijemo $L(M) = L(G)$. \qed

\begin{primjer}
KA za rijeci koje zavrsavaju na $ab$ iz poglavlja 3 pretvara se u DLG:
\begin{align*}
  V_{q_0} &\to a V_{q_a} \mid b V_{q_0} \\
  V_{q_a} &\to a V_{q_a} \mid b V_{q_{ab}} \\
  V_{q_{ab}} &\to a V_{q_a} \mid b V_{q_0} \mid \eps
\end{align*}
Ovo generira točno riječi koje završavaju na $ab$.
\end{primjer}

% ================================================================
\section{Gramaticki sustav i Ardenovo pravilo}
% ================================================================

Za smjer $\mathrm{DLG} \to \mathrm{RI}$ koristimo tehniku
\textbf{gramatickih jednadžbi}.

\subsection{Gramaticki sustav}

Svaka DLG zadaje \emph{sustav jednadžbi} nad jezicima: za svaku
varijablu $A$ skupimo sva njena pravila i zapisemo jednadžbu:
\[
  A = \alpha_1 B_1 \cup \alpha_2 B_2 \cup \cdots \cup \alpha_k B_k
  \;\cup\; (\eps \text{ ako ima pravilo } A \to \eps).
\]
Ako $A$ nema nijednog pravila, dodamo $A = \emptyset$.

Supstitucijom (uvrstavanjem jednadžbi jedne u drugu) pokusavamo
eliminirati varijable i doci do regularnog izraza za $S$.

\subsection{Problem: rekurzija}

Supstitucija funkcionira kada nema rekurzije. Ali DLG cesto ima
rekurzivna pravila poput $A = aA \cup b$, gdje se $A$ pojavljuje
s obje strane. Sto tada?

\begin{propozicija}[Ardenovo pravilo]
Jednadžba $X = AX \cup B$ ima rjesenje $X = A^* B$.
Ako je $A$ pozitivan, to je \emph{jedino} rjesenje.
Inace, to je \emph{najmanji} skup koji zadovoljava jednadžbu.
\end{propozicija}

\textit{Provjera.} Uvrstimo $X = A^*B$ u jednadžbu:
\[
  A(A^*B) \cup B = (AA^*)B \cup B = A^+B \cup \eps B = (A^+ \cup \eps)B = A^*B. \; \checkmark
\]

Ardenovo pravilo je analogon formule za geometrijski red:
$a \cdot r + b = r$ ima rjesenje $r = b/(1-a)$, tj.\ $r = a^* b$
u jezicnoj algebri.

\begin{primjer}[Pretvorba DLG $\to$ RI]
Pretvorimo DLG s pravilima:
\[
  S \to aA \mid bB, \quad
  A \to aA \mid bA \mid aN, \quad
  B \to aB \mid bB \mid bN, \quad
  N \to \eps.
\]
Gramaticki sustav ($N = \eps$ odmah uvrstimo):
\begin{align*}
  S &= aA \cup bB \\
  A &= aA \cup bA \cup a = (a \cup b)A \cup a
\end{align*}
Po Ardenovom pravilu ($a \cup b$ je pozitivan):
$A = (a \cup b)^* a$.

Analogno: $B = (a \cup b)^* b$.

Uvrstimo u $S$:
\[
  S = a(a \cup b)^*a \cup b(a \cup b)^*b.
\]
To su sve riječi koje počinju i završavaju istim slovom! \checkmark
\end{primjer}

% ================================================================
\section{RI $\to$ NKA: Thompsonova konstrukcija}
% ================================================================

Sad zatvaramo krug s druge strane: svaki regularni izraz mozemo
pretvoriti u NKA. Dokaz ide indukcijom po izgradnji izraza.

\begin{propozicija}
Svaki RI-regularni jezik je NKA-regularan.
\end{propozicija}

\textit{Dokaz (Thompsonova konstrukcija).}
Za svaki inicijalni jezik konstruiramo trivijalni NKA:
\begin{itemize}
  \item $\emptyset$: jedno stanje, nema prijelaza, nema zavrsnih stanja.
  \item $\eps$: jedno stanje, ono je i početno i završno, nema prijelaza.
  \item $\{\alpha\}$: dva stanja, prijelaz $\alpha$ od početnog do završnog.
\end{itemize}
Za slozenije izraze koristimo konstrukcije iz poglavlja 3
(unija s $\eps$-prijelazima, konkatenacija, zvijezda). \qed

% ================================================================
\section{Lijevolinearne gramatike i reverzi}
% ================================================================

\begin{definicija}[Lijevolinearna gramatika]
Gramatika je \textbf{lijevolinearna} (LLG) ako su sva njena pravila
oblika:
\[
  A \to B\alpha \qquad \text{ili} \qquad A \to \eps,
\]
dakle varijabla je na \emph{lijevom} kraju desne strane pravila.
\end{definicija}

LLG generiraju iste jezike kao i DLG (a time i sve ostale klase),
ali to nije odmah ocito. Dokazujemo to pomocu \textbf{reverza}.

\begin{definicija}[Reverz]
\textbf{Reverz} rijeci $w = \alpha_1 \alpha_2 \cdots \alpha_n$ je
$w^R = \alpha_n \cdots \alpha_2 \alpha_1$.
\textbf{Reverz jezika} je $L^R = \{w^R : w \in L\}$.
\end{definicija}

Za regularne izraze vrijede jednostavne formule:
\[
  \emptyset^R = \emptyset, \quad \eps^R = \eps, \quad \alpha^R = \alpha,
\]
\[
  (LM)^R = M^R L^R, \quad (L \cup M)^R = L^R \cup M^R, \quad (L^*)^R = (L^R)^*.
\]

Iz ovih formula indukcijom slijedi:

\begin{propozicija}
Reverz svakog RI-regularnog jezika je ponovo RI-regularan.
Dakle, klasa $\mathrm{Reg_{RI}}$ zatvorena je na reverz.
\end{propozicija}

Sada, DLG koja generira jezik $L$ lako se pretvori u LLG koja generira
$L^R$ (samo okrenimo redoslijed u pravilima). Kako je $\mathrm{Reg_{RI}}$
zatvorena na reverz, zakljucujemo:
\[
  \mathrm{Reg_{LLG}} = (\mathrm{Reg_{DLG}})^R = (\mathrm{Reg_{RI}})^R = \mathrm{Reg_{RI}}.
\]

% ================================================================
\section{Zakljucak: sve je isto}
% ================================================================

Evo kompletnog kruga pretvorbi:

\begin{center}
\renewcommand{\arraystretch}{1.6}
\begin{tabular}{lll}
\hline
\textbf{Od} & \textbf{Do} & \textbf{Metoda} \\
\hline
RI & NKA & Thompsonova konstrukcija (indukcija po izrazu) \\
NKA & KA & Partitivna konstrukcija \\
KA & DLG & Stanja $\leftrightarrow$ varijable, prijelazi $\leftrightarrow$ pravila \\
DLG & RI & Gramaticki sustav + Ardenovo pravilo \\
DLG & LLG & Reverz (okretanje pravila) \\
\hline
\end{tabular}
\end{center}

\begin{kljucno}
Sve ove pretvorbe su algoritamske --- dajte mi jedno, dat cu vam
bilo koje drugo. Svi formalizmi su jednako mocni; razlika je samo
u pogodnosti za razlicite svrhe. RI je dobar za opis, KA za
implementaciju, DLG za vezu s gramatikama.
\end{kljucno}

\begin{zadatak}
\begin{enumerate}[label=(\alph*), itemsep=4pt]
  \item Pretvorite regularni izraz $a^*(b \cup c)^+$ u NKA
        Thompsonovom konstrukcijom.
  \item Primijenite partitivnu konstrukciju na rezultat.
  \item Pretvorite KA koji ste dobili u DLG.
  \item Rijesite gramaticki sustav te DLG i dobijte regularni izraz.
        Usporedite s polaznim izrazom.
  \item Napisite LLG ekvivalentnu DLG iz (c).
\end{enumerate}
\end{zadatak}

% ================================================================
\section*{Sazetak}
% ================================================================

Pokazali smo da pet naizgled razlicitih formalizma opisuju istu klasu:

\medskip
\noindent
\begin{tabular}{ll}
  $\mathrm{Reg_{KA}}$ & jezici KA \\
  $\mathrm{Reg_{NKA}}$ & jezici NKA \\
  $\mathrm{Reg_{RI}}$ & regularni izrazi \\
  $\mathrm{Reg_{DLG}}$ & desnolinearne gramatike \\
  $\mathrm{Reg_{LLG}}$ & lijevolinearne gramatike \\
\end{tabular}

\medskip
Sve su iste klase. Tu klasu zovemo $\mathrm{Reg}$ --- \textbf{regularni
jezici}. U sljedecem poglavlju prelazimo na bezkontekstne jezike ---
prvu pravu nadklasu.


% Chapter 5 — includable by main.tex
\chapter{Bezkontekstni jezici}
\label{chap:5}

\begin{quote}
\itshape
Regularni jezici ne mogu prebrojati. Čim je potrebno "zapamtiti"
proizvoljan broj i usporediti ga s necim drugim, regularnost otkazuje.
Bezkontekstni jezici rješavaju upravo taj problem --- dodavanjem stoga.
Ovo poglavlje uvodi bezkontekstne gramatike, potisne automate,
dokazuje njihovu ekvivalentnost, i opisuje ključni algoritam CYK.
\end{quote}


% ================================================================
\section{Motivacija i definicija BKG}
% ================================================================

Prisjetimo se: jezik $L_= = \{a^n b^n \mid n \geq 1\}$ nije regularan
jer KA ne može istovremeno pratiti broj $a$-ova i usporedivati ga s
brojem $b$-ova --- konačno mnogo stanja nije dovoljno.

No postoji gramatika koja ga generira:
\[
  S \to aSb \mid ab.
\]
Lako se vidi indukcijom da ova gramatika generira točno $a^n b^n$:
pravilo $S \to aSb$ dodaje po jedan $a$ s lijeva i jedan $b$ s desna,
a $S \to ab$ zatvara rekurziju.

Ova gramatika je \textbf{bezkontekstna}: svako pravilo zamjenjuje
jednu varijablu nizom simbola, bez obzira na kontekst oko nje.

\begin{definicija}[Bezkontekstna gramatika]
\textbf{Bezkontekstna gramatika} (BKG) je gramatika $G = (V, \Sig, \to, S)$
u kojoj je svako pravilo oblika
\[
  A \to w, \qquad A \in V,\; w \in (\Sig \cup V)^*.
\]
Dakle, lijeva strana je uvijek točno jedna varijabla. Jezik koji
$G$ generira je $L(G) = \{w \in \Sig^* : S \Rightarrow^* w\}$.
Jezik je \textbf{bezkontekstan} (BK) ako ga generira neka BKG.
\end{definicija}

\begin{napomena}
Svaka DLG je BKG (pravila $A \to \alpha B$ i $A \to \eps$ očito
su bezkontekstna). Dakle, $\mathrm{Reg} \subseteq \mathrm{BK}$.
Jezik $L_=$ pokazuje da je inkluzija prava: $\mathrm{Reg} \subsetneq \mathrm{BK}$.
\end{napomena}

\begin{primjer}[Aritmeticki izrazi]
BKG za jednostavne aritmeticke izraze nad varijablom $x$ i brojevima:
\begin{align*}
  E &\to E + T \mid T \\
  T &\to T * F \mid F \\
  F &\to (E) \mid x \mid 0 \mid 1 \mid \ldots
\end{align*}
Ova gramatika odrazava prioritet operatora: mnozenje ($T$) ima visi
prioritet od zbrajanja ($E$), a zagrade ($F$) imaju najvisi.
Iz $E$ možemo izvesti npr.\ $x + x * x$, i stablo parsiranja
automatski reflektira ispravnu strukturu $(x + (x * x))$.
\end{primjer}

\begin{primjer}[Zadatak sa slajdova]
Koji jezik generira gramatika:
\[
  S \to 0J \mid 1N \mid \eps, \quad J \to 1S \mid SJ, \quad N \to 0S \mid SN\,?
\]
(Svaki put kad iz $S$ odemo u $0J$, sljedeći znak mora biti $1$)
(jer $J \to 1S$), sto znaci svaki $0$ ima par $1$ iza sebe.
Slicno, svaki $1$ ima par $0$ iza sebe. Ovo generira jezik svih
binarnih rijeci s jednakim brojem $0$-ova i $1$-ova --- i taj je
jezik bezkontekstan, ali nije regularan.
\end{primjer}

% ================================================================
\section{Zatvorenost klase BK}
% ================================================================

\begin{propozicija}
Klasa $\mathrm{BK}$ zatvorena je na uniju, konkatenaciju i Kleeneve operacije.
\end{propozicija}

\textit{Dokaz (skica).}
Za BKG $G_1$ i $G_2$ s disjunktnim skupovima varijabli koje generiraju
$L_1$ i $L_2$:
\begin{itemize}
    \item \textbf{Unija}: dodamo novu početnu varijablu $S$ i pravilo
      $S \to S_1 \mid S_2$.
  \item \textbf{Konkatenacija}: dodamo $S$ i pravilo $S \to S_1 S_2$.
  \item \textbf{Zvijezda}: dodamo $S$ i pravilo $S \to S_1 S \mid \eps$.
\end{itemize}

\begin{kljucno}
Za razliku od regularnih jezika, klasa BK \textbf{nije} zatvorena na presjek
ni na komplement. Ovo je bitna razlika: $\mathrm{Reg}$ je zatvorena na sve
skupovne operacije, a $\mathrm{BK}$ nije. Dokaz dolazi u poglavlju 6
(lema o pumpanju za BK).

Ipak vrijedi \textbf{konusno svojstvo}: presjek bezkontekstnog i
regularnog jezika je bezkontekstan.
\end{kljucno}

% ================================================================
\section{Potisni automati (PA)}
% ================================================================

KA je premoczan za regularne, a premoczan --- ne, preslabacak ---
za bezkontekstne jezike. Treba mu memorija. Najjednostavniji dodataak:
\textbf{stog} (stack, magacin).

\begin{definicija}[Potisni automat]
\textbf{Potisni automat} (PA) je uredena sestorka
\[
  M = (Q,\, \Sig,\, \Gamma,\, \Delta,\, q_0,\, F),
\]
gdje su $Q$, $\Sig$, $q_0$, $F$ kao kod NKA, a novo je:
\begin{itemize}
  \item $\Gamma$ --- \textbf{abeceda stoga} (može se razlikovati od $\Sig$),
  \item $\Delta \subseteq Q \times \Sig_\eps \times \Gamma_\eps \times Q \times \Gamma_\eps$
        --- relacija prijelaza.
\end{itemize}
Prijelaz $(q, \alpha, t, p, s) \in \Delta$ znači: ``u stanju $q$, čitajući
$\alpha$ (ili $\eps$), s $t$ na vrhu stoga (ili $\eps$ ako stog ne gledamo),
pređi u stanje $p$ i zamijeni $t$ s $s$ na stogu.''
\end{definicija}

Prijelaz može kombinirati različite akcije:
\begin{center}
\renewcommand{\arraystretch}{1.4}
\begin{tabular}{lll}
\hline
$t$ & $s$ & Sto se dogada \\
\hline
$\eps$ & $\gamma$ & \textbf{push}: stavi $\gamma$ na stog \\
$\gamma$ & $\eps$ & \textbf{pop}: skini $\gamma$ sa stoga \\
$\gamma$ & $\gamma$ & provjeri vrh stoga (ostavi nepromijenjenim) \\
$\eps$ & $\eps$ & ignoriraj stog \\
\hline
\end{tabular}
\end{center}

\begin{napomena}
Za PA je \textbf{nedeterminizam ključan}. Deterministički PA prepoznaju
uzu klasu od svih BK-jezika (postoje BK-jezici koje deterministički
PA ne može prepoznati). Ovo je bitna razlika od KA, gdje
determinizam i nedeterminizam daju istu klasu.
\end{napomena}

\subsection{Jednostavni PA (JPA)}

Radi lakse teorije, definiramo \textbf{jednostavni PA} (JPA) koji ima
tri dodatna svojstva:
\begin{enumerate}
  \item Ima točno jedno završno stanje $q_X$.
  \item Kad dosegne $q_X$, stog mu je prazan.
  \item Svaki prijelaz je isključivo \emph{push} ili \emph{pop} (ne oboje).
\end{enumerate}

Svaki PA može se pretvoriti u ekvivalentni JPA: uvedemo novi simbol
$\$$ na dno stoga, sva bivsa zavrsna stanja isprazne stog i odu u $q_X$,
a svaki ``kombinirani'' prijelaz rastavimo na dva.

\begin{primjer}[PA za $L_=$]
Konstruirajmo PA koji prepoznaje $L_= = \{a^n b^n \mid n \geq 1\}$.

Ideja: stavljamo $a$-ove na stog dok citamo $a$-ove, pa skidamo
po jedan za svaki $b$ koji procitamo.

\begin{itemize}
  \item Stanja: $\{q_0, q_1, q_\text{acc}\}$, početno $q_0$, završno $\{q_\text{acc}\}$. 
  \item Na početku push $\$$ (marker dna stoga).
  \item U $q_0$: za svaki $a$, push $a$ (ostani u $q_0$).
  \item Prijelaz u $q_1$: $\eps$-prijelaz (početak $b$-ova).
  \item U $q_1$: za svaki $b$, pop $a$.
  \item Kad je na vrhu stoga $\$$, prijelaz u $q_\text{acc}$.
\end{itemize}
\end{primjer}

% ================================================================
\section{Ekvivalentnost BKG i PA}
% ================================================================

\begin{propozicija}[Sipser, lema 2.21 i teorem 2.20]
Jezik je bezkontekstan ako i samo ako ga prepoznaje neki PA.
\end{propozicija}

\subsection{BKG $\to$ PA}

Konstruiramo PA koji ``pogada'' izvod riječi iz BKG:

\begin{enumerate}
  \item Push $\$$ i početnu varijablu $S$ na stog.
  \item Ponavljaj:
    \begin{itemize}
      \item Ako je na vrhu stoga varijabla $A$: nedeterministički
        odaberi pravilo $A \to w$ i zamijeni $A$ sa simbolima od $w$
        (u obrnutom redoslijedu, da prvi simbol bude na vrhu).
      \item Ako je na vrhu stoga znak $\alpha \in \Sig$: pročitaj
        sljedeći ulazni znak; ako se slaže, pop; inače odbaci granu.
      \item Ako je na vrhu stoga $\$$: prihvati.
    \end{itemize}
\end{enumerate}

Ključna napomena: ovaj PA koristi nedeterminizam da ``pogodi'' točni
lijevi izvod. PA prolazi točno kroz konfiguracije koje odgovaraju
koracima lijevog izvoda.

\subsection{PA $\to$ BKG}

Ovaj smjer je tehnički zahtjevniji. Svaki PA pretvorimo u JPA,
a zatim za svaki par stanja $(p, q)$ uvedemo varijablu ${}_{p}V_q`
sa značenjem: ``JPA iz stanja $p$ s praznim stogom može pročitati
neku riječ i doći u stanje $q$ s praznim stogom''.

Pravila gramatike odrazavaju dva slucaja:
\begin{itemize}
    \item \textbf{Stog postane prazan u međuvremenu}: postoji stanje $r$
      gdje stog postane prazan, pa se izvod ``lomi'' na dva dijela:
        \[ {}_{p}V_q \to {}_{p}V_r \; {}_{r}V_q. \]
  \item \textbf{Stog nikad nije prazan}: isti simbol $\gamma$ koji
        je pocetak push na kraju pop-a. Ako je push na pocetku
        $(p, \alpha, \eps, r, \gamma)$ i pop na kraju $(s, \beta, \gamma, q, \eps)$:
        \[ {}_{p}V_q \to \alpha \; {}_{r}V_s \; \beta. \]
\end{itemize}
Pocetna varijabla je ${}_{q_0}V_{q_X}$.

% ================================================================
\section{Chomskyjeva normalna forma (ChNF)}
% ================================================================

Svaka BKG može se pretvoriti u kanonski oblik koji je koristan
za algoritme i dokaze.

\begin{definicija}[Chomskyjeva normalna forma]
BKG je u \textbf{Chomskyjevoj normalnoj formi} (ChNF) ako su sva
njena pravila jednog od dva oblika:
\[
  A \to BC \qquad \text{(sirenje)}, \qquad A \to \alpha \qquad \text{(citanje)},
\]
gdje su $A, B, C \in V$ varijable (s uvjetom $S \notin \{B, C\}$),
a $\alpha \in \Sig$ znak.
\end{definicija}

Svaka ChNF gramatika generira pozitivan jezik (jer ni sirenje ni
citanje ne mogu dati $\eps$). Dokazat cemo i obrat: svaki pozitivni
BK-jezik ima ChNF gramatiku.

\subsection{Algoritam pretvorbe u ChNF (5 faza)}

Pretvorba ide u pet faza, svaka osigurava jedno svojstvo ne
kvareci prethodna.

\medskip
\noindent\textbf{Faza START.}
Ako se pocetna varijabla $S$ pojavljuje na desnoj strani ijednog
pravila, uvedemo novu varijablu $S_0$ s pravilom $S_0 \to S$ i
proglasimo $S_0$ novom pocetnom varijablom.

\medskip
\noindent\textbf{Faza TERM.}
Za svaki znak $\alpha \in \Sig$ koji se pojavljuje u nekom pravilu
uz jos nesto (dakle nije sam), uvedemo novu varijablu $N_\alpha$
s pravilom $N_\alpha \to \alpha$, i sve takve pojave $\alpha$
zamijenimo s $N_\alpha$.

\emph{Primjer:} pravilo $A \to aBc$ postaje $A \to N_a B N_c$
(gdje su $N_a \to a$ i $N_c \to c$ nova pravila citanja).

\medskip
\noindent\textbf{Faza BIN.}
Za svako pravilo s $n \geq 3$ simbola na desnoj strani,
$A \to V_1 V_2 \cdots V_n$, uvedemo nove varijable $A_1, \ldots, A_{n-2}$
i rastavimo ga na lanac:
\[
  A \to V_1 A_1, \quad A_1 \to V_2 A_2, \quad \ldots, \quad A_{n-2} \to V_{n-1} V_n.
\]

\medskip
\noindent\textbf{Faza DEL.}
Eliminiramo pravila oblika $A \to \eps$ (za $A \neq S$).
Algoritam:
\begin{enumerate}
  \item Nademo sve \textbf{nepozitivne varijable} $\mathrm{np} = \{A : A \Rightarrow^* \eps\}$
        iteracijom: pocnemo s $\mathrm{np} = \emptyset$, i dodajemo $A$
        kad su sve varijable na desnoj strani nekog pravila za $A$ u $\mathrm{np}$.
  \item Za svako pravilo koje sadrzi $k > 0$ varijabli iz $\mathrm{np}$,
        dodamo $2^k$ novih pravila (jedno za svaki podskup tih pojava
        koje brisemo), i brisemo originalno pravilo $A \to \eps$.
\end{enumerate}
Gramatika je bila pozitivna ako i samo ako $S \notin \mathrm{np}$.

\medskip
\noindent\textbf{Faza UNIT.}
Eliminiramo pravila oblika $A \to B$ (jedna varijabla na desnoj strani).
Za svako takvo pravilo: uklonimo ga, zapamtimo, i za svako pravilo
$B \to u$ dodamo $A \to u$ (osim ako smo to pravilo vec uklonili).

\begin{primjer}[Pretvorba u ChNF]
Pretvorimo gramatiku $S \to ASA \mid aB$, $A \to B \mid S$, $B \to b \mid \eps$.

\textbf{START}: $S$ se ne pojavljuje na desnoj strani --- preskacemo.

\textbf{TERM}: $a$ i $b$ se pojavljuju uz ostalo: uvedemo $N_a \to a$, $N_b \to b$.
Dobijemo: $S \to ASA \mid N_a B$, $B \to N_b \mid \eps$.

\textbf{BIN}: $S \to ASA$ ima tri simbola: $S \to A A_1$, $A_1 \to SA$.

\textbf{DEL}: $B \to \eps$, pa $\mathrm{np} = \{B\}$.
Pravilo $S \to N_a B$ daje i $S \to N_a$.

\textbf{UNIT}: $A \to B$ i $A \to S$. Za $A \to B$: $B$ ima pravila
$B \to N_b$, pa dodamo $A \to N_b$. Za $A \to S$: $S$ ima pravila
$S \to A A_1 \mid N_a B \mid N_a$, pa dodamo ta pravila za $A$.
\end{primjer}

% ================================================================
\section{Algoritam CYK}
% ================================================================

ChNF nam omogucuje efikasan algoritam parsiranja.

\begin{propozicija}
Problem prihvaćanja $A_\mathrm{BK}$ je odlučiv. Konkretno,
algoritam CYK rješava ga u $O(n^3 \cdot |G|)$ vremenu,
gdje je $n = |w|$ duljina ulazne riječi.
\end{propozicija}

\subsection{Ideja: dinamičko programiranje}

Ključna opservacija za ChNF: ako $A \Rightarrow^* w_{i:j}$
(podriječ od pozicije $i$ do $j$), tada prvi korak mora biti širenje
$A \to BC$, a $B$ i $C$ izvode neprazne podrijeci. Dakle postoji
$k \in (i, j)$ takav da $B \Rightarrow^* w_{i:k}$ i $C \Rightarrow^* w_{k:j}$.

$Ovo sugerira \textbf{dinamičko programiranje}: rješavamo podprobleme
$A \mathop{?} w_{i:j}$ redom od kracih prema duljima.

\subsection{Tablica i algoritam}

Tablicu $T$ dimenzija $|V| \times n \times n$ popunjavamo:

\medskip
\noindent\textbf{Baza} ($j = i + 1$, jednoznakovna podriječ):
\[
  T[A][i][i+1] = \mathrm{True} \iff (A \to w_i) \in G.
\]

\noindent\textbf{Korak} ($j - i = \ell > 1$):
\[
  T[A][i][j] = \mathrm{True} \iff \exists\, (A \to BC) \in G,\; k \in (i,j) :
  T[B][i][k] \wedge T[C][k][j].
\]

\noindent\textbf{Odgovor}: $w \in L(G) \iff T[S][0][n] = \mathrm{True}$.

\begin{primjer}[CYK na djelu]
Neka je gramatika (u ChNF): $S \to AB \mid BC$, $A \to BA \mid a$,
 $B \to CC \mid b$, $C \to AB \mid a$ i riječ $w = baaba$.

Duljina je $5$, tablica je $5 \times 5$. Popunjavamo redak po redak
(duljine $1$, $2$, $3$, $4$, $5$):

\begin{center}
\renewcommand{\arraystretch}{1.3}
\begin{tabular}{c|ccccc}
  & $b$ & $a$ & $a$ & $b$ & $a$ \\
  & $w_{0:1}$ & $w_{1:2}$ & $w_{2:3}$ & $w_{3:4}$ & $w_{4:5}$ \\\hline
  $S$ & & $\{S\}$ ok? & \ldots & \ldots & \\
  $A$ & & $A$ & $A,C$ & & $A,C$ \\
  $B$ & $B$ & & & $B$ & \\
  $C$ & & $A,C$ & $A,C$ & & $A,C$ \\
\end{tabular}
\end{center}

(Za punu tablicu pogledajte Sipser poglavlje 2.3 ili
\url{https://youtu.be/EloKZv84hUc?t=246}.)

Na kraju provjerimo $T[S][0][5]$ --- ako je True, riječ je u jeziku.
\end{primjer}

\subsection{Rekonstrukcija stabla parsiranja}

Osim bool vrijednosti, u tablicu možemo zapisati i ``svjedoke'':
trojke $(B, C, k)$ koje su dovele do True. Tada možemo rekonstruirati
cijelo stablo parsiranja rekurzivnim obilaskom tablice.

\begin{kljucno}
CYK je kubne \emph{vremenske} i kvadratne \emph{prostorne} složenosti
za jednoznačne gramatike. Nedostatak: radi na ChNF, ne na originalnoj
gramatici, pa stablo parsiranja sadrži ``artifakte'' pretvorbe u ChNF.
\end{kljucno}

% ================================================================
\section{Stabla parsiranja i viseznacanost}
% ================================================================

\begin{definicija}[Stablo parsiranja]
\textbf{Stablo parsiranja} iz BKG $G$ za riječ $w$ je stablo u cijem
korijenu je pocetna varijabla, svaki unutarnji cvor je varijabla
cija djeca odgovaraju desnoj strani nekog pravila, a listovi cine
riječ $w$ (čitani slijeva nadesno).
\end{definicija}

Vrijedi: $w \in L(G)$ ako i samo ako postoji stablo parsiranja za $w$.

\begin{definicija}[Viseznacanost]
Riječ $w$ je \textbf{$G$-visezna\-cna} ako postoji više različitih
stabala parsiranja za $w$ u $G$. Gramatika je \textbf{visezna\-cna}
ako postoji visezna\-cna riječ. Jezik je \textbf{visezna\-can} ako
je svaka njegova BKG visezna\-cna.
\end{definicija}

\begin{primjer}
Gramatika $S \to SS \mid a$ je visezna\-cna: riječ $aaa$ ima dva
razlicita stabla parsiranja ($(aa)a$ i $a(aa)$).

Pazi: vise izvoda ne znaci nuzno visezna\-cnost! Bitna su stabla
parsiranja, ne izvodi. Lijevi izvod je jedinstven po stablu parsiranja
--- postoji bijekcija lijevi izvodi $\leftrightarrow$ stabla parsiranja.
\end{primjer}

% ================================================================
\section*{Sazetak}
% ================================================================

\medskip\noindent
\begin{tabular}{ll}
  BKG & bezkontekstna gramatika: $A \to w$ \\
  PA & potisni automat s stogom \\
  JPA & jednostavni PA (jedan ulaz, push/pop odvojeni) \\
  ChNF & $A \to BC$ ili $A \to \alpha$; korisna za algoritme \\
  CYK & $O(n^3)$ algoritam parsiranja pomoću dinamičkog programiranja \\
  Viseznacanost & više stabala parsiranja za istu riječ \\
\end{tabular}

\medskip\noindent
U sljedećem poglavlju dokazujemo da neke jezike nije moguće opisati BKG-om
--- pomoću \textbf{leme o pumpanju za BK}.


% Chapter 6 — includable by main.tex
\chapter{Leme o pumpanju}
\label{chap:6}

\begin{quote}
\itshape
Ta stvar zove se \textbf{svojstvo pumpanja}.

Intuicija kaže "treba pamtiti neograničeno puno", ali to nije formalni dokaz.
Leme o pumpanju daju nam precizno oružje: ako jezik nema određeno
"periodično" svojstvo, onda ne može biti u toj klasi.
\end{quote}

\section{Lema o pumpanju za regularne jezike}
% ================================================================

\begin{propozicija}[Lema o pumpanju, Sipser 1.70]
Neka je $L$ regularan jezik. Tada postoji broj $p \geq 1$
(``duljina pumpanja'') takav da se svaka riječ $w \in L$ duljine
$|w| \geq p$ može zapisati kao $w = xyz$ uz uvjete:
\begin{enumerate}[label=(\roman*)]
  \item $y \neq \eps$ \quad (pumpa je neprazna),
  \item $|xy| \leq p$ \quad (pumpa je blizu početka),
  \item $xy^i z \in L$ za sve $i \geq 0$ \quad (pumpanje cuva jezik).
\end{enumerate}
\end{propozicija}

\textit{Dokaz.}
Neka KA $M$ prepoznaje $L$, i neka je $p = |Q|$ broj stanja.
Za riječ $w \in L$ s $|w| \geq p$: čitajući prvih $p$ znakova,
automat prolazi kroz $p+1$ stanja $r_0, r_1, \ldots, r_p$.
Po \textbf{Dirichletovom principu} (golubinjak), neka dva stanja
su ista: $r_i = r_j$ za neke $0 \leq i < j \leq p$.

Definirajmo:
\[
  x = w_{0:i}, \quad y = w_{i:j}, \quad z = w_{j:|w|}.
\]
Tada:
\begin{itemize}
  \item $y \neq \eps$ jer $j > i$, dakle $|y| = j - i > 0$.
  \item $|xy| = j \leq p$.
  \item Automat koji procita $xy^i z$: ulazi u $r_i$ nakon $x$,
        petlja $i$ puta kroz $r_i \to \cdots \to r_j = r_i$,
        zatim nastavlja od $r_j = r_i$ kroz $z$ i zavrsava u istom
        zavrsnom stanju kao i za $w$. \qed
\end{itemize}

\begin{kljucno}
Lemu koristimo \textbf{kontrapositivno}: ako za svaki $p$ možemo
naći riječ $w \in L$ koja se ne može pumpati (za svaki rastav
$xyz$ koji zadovoljava (i) i (ii), postoji $i$ takav da
$xy^iz \notin L$), tada $L$ nije regularan.

Dokaz je \textbf{igra}: naš protivnik odabire $p$, mi biramo $w$,
protivnik bira rastav $xyz$, mi biramo $i$. Pobjedujemo ako dobijemo
$xy^iz \notin L$.
\end{kljucno}

% ================================================================
\section{Primjena: $L_=$ nije regularan}
% ================================================================

\begin{propozicija}
Jezik $L_= = \{a^n b^n \mid n \geq 1\}$ nije regularan.
\end{propozicija}

\textit{Dokaz.}
Pretpostavimo suprotno: neka je $L_=$ regularan s duljinom pumpanja $p$.

Odaberemo riječ $w = a^p b^p \in L_=$ (duljine $2p \geq p$, dakle
zadovoljava uvjet). Prema lemi, postoji rastav $w = xyz$ s
$y \neq \eps$ i $|xy| \leq p$.

Iz $|xy| \leq p = |a^p|$ slijedi da $x$ i $y$ leze unutar bloka $a$-ova,
dakle $y = a^m$ za neki $m \geq 1$.

Pumpajmo s $i = 0$: $xy^0z = xz = a^{p-m}b^p$. Da bi ovo bilo u $L_=$,
trebamo jednak broj $a$-ova i $b$-ova: $p - m = p$, ali $m \geq 1$
pa $p - m < p$. Kontradikcija: $xz \notin L_=$. \qed

\begin{primjer}[Tipična strategija dokaza]
Struktura dokaza je uvijek ista:
\begin{enumerate}
  \item \textbf{Pretpostavimo} da je $L$ regularan s duljinom pumpanja $p$.
    \item \textbf{Odaberemo} konkretnu riječ $w \in L$ duljine $\geq p$
      (tipično $w$ ovisi o $p$, npr.\ $w = a^p b^p$).
    \item \textbf{Analiziramo} sve moguće rastave $xyz$ koji zadovoljavaju
      uvjete (i) i (ii). Često uvjet $|xy| \leq p$ snažno ograničava
      gdje može ležati $y$.
  \item \textbf{Pumpamo} s nekim $i$ (tipicno $i = 0$ ili $i = 2$) i
        pokazujemo da $xy^i z \notin L$.
  \item \textbf{Kontradikcija}: $L$ nije regularan.
\end{enumerate}
\end{primjer}

\begin{primjer}[$\{w \in \{a,b\}^* : |w|_a = |w|_b\}$ nije regularan]
Neka je $L = \{w : |w|_a = |w|_b\}$. Odaberemo $w = a^p b^p$.
Sve kao gore: $y = a^m$, pumpanje s $i = 0$ daje $a^{p-m}b^p$
s razlicitim brojem $a$-ova i $b$-ova. \qed
\end{primjer}

\begin{primjer}[$\{0^n : n \text{ je prost}\}$ nije regularan]
Neka je $L = \{0^n : n \text{ je prost}\}$.
Odaberemo $w = 0^p$ gdje je $p$ prvi prost broj $\geq p$ duljine
pumpanja (takav postoji jer ima beskonačno prostih).

Rastav: $w = xyz$, $|y| = m \geq 1$, $|xy| \leq p$.
Pumpamo s $i = |xz| + 1$: duljina od $xy^i z$ je
$|x| + i \cdot m + |z| = |xz| + im = |xz|(1 + m) + m \cdot |xz|$\ldots

Koristimo drugi trik: $|xy^i z| = |xz| + im$.
Uzmemo $i = |xz|$, tada $|xy^iz| = |xz| + |xz| \cdot m = |xz|(1+m)$
--- djeljivo s $|xz|$ i s $(1+m)$, dakle nije prost za $|xz| > 1$.
\qed
\end{primjer}

\begin{zadatak}
Dokazite da sljedeći jezici nisu regularni (koristeći lemu o pumpanju):
\begin{enumerate}[label=(\alph*)]
  \item $\{a^n b^{2n} \mid n \geq 0\}$
  \item $\{w \in \{a,b\}^* : w = w^R\}$ (palindromi)
  \item $\{1^{n^2} \mid n \geq 0\}$ (kvadrati u unarnom zapisu)
  \item $\{a^i b^j c^k \mid i + j = k\}$
\end{enumerate}
\end{zadatak}

% ================================================================
\section{Lema o pumpanju za bezkontekstne jezike}
% ================================================================

Ista ideja, ali za BK-jezike i s pet dijelova umjesto tri.

\begin{propozicija}[Lema o pumpanju za BK, Sipser 2.34]
Neka je $L$ bezkontekstan jezik. Tada postoji broj $p \geq 1$
takav da se svaka riječ $s \in L$ s $|s| \geq p$ može zapisati
kao $s = uvxyz$ uz uvjete:
\begin{enumerate}[label=(\roman*)]
  \item $vy \neq \eps$ \quad (barem jedna od pumpi je neprazna),
  \item $|vxy| \leq p$ \quad (pumpe su blizu jedna drugoj),
  \item $uv^i x y^i z \in L$ za sve $i \geq 0$ \quad (pumpanje cuva jezik).
\end{enumerate}
\end{propozicija}

\textit{Ideja dokaza.}
Neka BKG $G$ (u ChNF) generira $L$, i neka je $b = \max |desna strana|`
(za ChNF, $b = 2$). Za dovoljno dugu riječ $s$, stablo parsiranja
mora imati duboku granu --- toliko duboku da se neka varijabla $R$
mora ponoviti duž iste grane (po Dirichletovom principu, ako
grana ima više od $|V|$ čvorova).

Istaknemo zadnja dva pojavljivanja $R$ u toj grani:
\[
  S \Rightarrow^* uRz \Rightarrow^* uvRyz \Rightarrow^* uvxyz = s.
\]
Budući da $R \Rightarrow^* vRy$, možemo iterirati: $R \Rightarrow^* v^i R y^i$,
a zatim zatvoriti s $R \Rightarrow^* x$. Time dobijemo $uv^i xy^i z \in L$
za svaki $i \geq 0$.

Uvjet $|vxy| \leq p$ dolazi iz toga da biramo \emph{zadnja} dva
pojavljivanja $R$ u dubini: podstablo ispod nižeg $R$ nije predugo,
pa riječ koju ono generira nije preduga. \qed

% ================================================================
\section{Primjena: $L_\equiv$ nije bezkontekstan}
% ================================================================

\begin{propozicija}
Jezik $L_\equiv = \{a^n b^n c^n \mid n \geq 1\}$ nije bezkontekstan.
\end{propozicija}

\textit{Dokaz.}
Pretpostavimo suprotno: neka je $p$ duljina pumpanja za $L_\equiv$.

Odaberemo $s = a^p b^p c^p$ (duljine $3p \geq p$). Prema lemi,
postoji rastav $s = uvxyz$ s $vy \neq \eps$ i $|vxy| \leq p$.

Ključna opservacija: $|vxy| \leq p$, a $a$-ovi i $c$-ovi su udaljeni
$2p > p$ znakova. Dakle $vxy$ ne može sadržavati i $a$-ove i $c$-ove
istovremeno.

\textbf{Slučaj 1}: $vxy$ ne sadrzi $c$-ove.
Tada je svih $p$ slova $c$ u dijelu $z$. Pumpamo s $i = 2$:
riječ $uv^2xy^2z$ ima više $a$-ova ili $b$-ova (jer smo pumpali $v$
ili $y$ koji sadrže $a$-ove ili $b$-ove), ali i dalje $p$ slova $c$.
Dakle $uv^2xy^2z \notin L_\equiv$.

\textbf{Slučaj 2}: $vxy$ ne sadrzi $a$-ove (analogno, ili BNSO).
Svih $p$ slova $a$ je u $u$. Pumpamo s $i = 0$:
$uxz$ ima $p$ slova $a$, ali manje od $p$ slova $b$ ili $c$
(jer $vy \neq \eps$ pa smo nešto uklonili). Dakle $uxz \notin L_\equiv$. \qed

\begin{primjer}[Nezatvorenost BK na presjek]
Definiramo:
\begin{align*}
  L_{a=b} &= \{a^n b^n c^k : n, k \geq 1\} = L_= \cdot c^+, \\
  L_{b=c} &= \{a^k b^n c^n : n, k \geq 1\} = a^+ \cdot L_=.
\end{align*}
Oba su bezkontekstna (konkatenacija BK i regularnog).
Ali $L_{a=b} \cap L_{b=c} = L_\equiv$ --- a upravo smo dokazali da
$L_\equiv$ nije bezkontekstan!

Zaključak: \textbf{BK nije zatvorena na presjek}. Iz de Morgana slijedi
ni na komplement.
\end{primjer}

\begin{zadatak}
Dokazite da sljedeći jezici nisu bezkontekstni:
\begin{enumerate}[label=(\alph*)]
  \item $\{a^n b^n c^n d^n \mid n \geq 1\}$
  \item $\{ww \mid w \in \{a,b\}^*\}$ (kvadrati)
  \item $\{a^i b^j c^k \mid i < j < k\}$
  \item $\{1^{n^2} \mid n \geq 0\}$ (je li ovo BK?)
\end{enumerate}
\end{zadatak}

% ================================================================
\section{Usporedba dviju lema}
% ================================================================

\begin{center}
\renewcommand{\arraystretch}{1.6}
\begin{tabular}{lll}
\hline
& \textbf{Reg (lema za Reg)} & \textbf{BK (lema za BK)} \\
\hline
Rastav & $xyz$ (3 dijela) & $uvxyz$ (5 dijelova) \\
Pumpe & $y$ (jedna pumpa) & $v$ i $y$ (dvije pumpe) \\
Uvjet nepraznosti & $y \neq \eps$ & $vy \neq \eps$ \\
Uvjet lokalizacije & $|xy| \leq p$ & $|vxy| \leq p$ \\
Pumpanje & $xy^iz \in L$ & $uv^ixy^iz \in L$ \\
Ključni alat dokaza & Dirichlet na stanjima KA & Dirichlet na varijablama u stablu \\
\hline
\end{tabular}
\end{center}

Obje leme imaju isti oblik dokaza neregularnosti/nebezkontekstnosti:
pretpostavi, odaberi rijec, analiziraj rastav, pumpa, kontradikcija.

\begin{kljucno}
Leme o pumpanju su \textbf{nuzni} uvjeti za regularnost/bezkontekstnost,
ali ne i \textbf{dovoljni}! Postoje jezici koji zadovoljavaju lemu o
pumpanju za Reg, a nisu regularni. Lema se koristi samo za dokazivanje
\emph{neregularnosti} --- ne može dokazati regularnost.
\end{kljucno}

% ================================================================
\section*{Sazetak}
% ================================================================

\medskip\noindent
\begin{tabular}{ll}
  Lema za Reg & $w = xyz$, $y \neq \eps$, $|xy| \leq p$, $xy^iz \in L$ \\
  Lema za BK  & $s = uvxyz$, $vy \neq \eps$, $|vxy| \leq p$, $uv^ixy^iz \in L$ \\
  Strategija & pretpostavi $\to$ odaberi riječ $\to$ analiziraj $\to$ pumpa $\to$ kontradikcija \\
\end{tabular}

\medskip\noindent
U sljedećem poglavlju prelazimo na kontekstne jezike i ograničene
automate --- prvu klasu za kojom trebamo ``pravu'' traku.


% Chapter 7 — includable by main.tex
\chapter{Kontekstni jezici}
\label{chap:7}

\begin{quote}
\itshape
Bezkontekstni jezici ne mogu prebrojati tri stvari odjednom.
Kontekstni jezici rješavaju i taj problem --- dodavanjem trake
ograničene na duljinu ulaza. Klasa Kon je u mnogo čemu lijepa:
zatvorena je na sve skupovne operacije, a prepoznavanje je odlučivo.
\end{quote}


% ================================================================
\section{Motivacija: što ne može BK?}
% ================================================================

U poglavlju 6 dokazali smo da $L_\equiv = \{a^n b^n c^n \mid n \geq 1\}$
nije bezkontekstan. Intuitivno: potisni automat može uskladiti
broj $a$-ova i $b$-ova (stavi $a$-ove na stog, skidaj za svaki $b$),
ali jednom kad skine sve s hrpe, više ne zna koliko je $b$-ova prebrojio ---
to je točno potrebno za provjeru $c$-ova.

Rješenje: umjesto stoga koji nestaje, trebamo traku po kojoj se
možemo kretati u oba smjera. Ključno ograničenje: traka je dugačka
\emph{točno koliko i ulazna riječ}. Nema više memorije od ulaza.

% ================================================================
\section{Ograniceni automat (OA)}
% ================================================================

\begin{definicija}[Ograniceni automat]
\textbf{Ograniceni automat} (OA) je uredena sestorka
\[
  O = (Q,\, \Sig,\, \Gamma,\, \Delta,\, q_0,\, q_X),
\]
gdje su $Q$, $\Sig \subseteq \Gamma$, $q_0$, $q_X$ kao obicno, a
\[
  \Delta \subseteq Q \times \Gamma \times Q \times \Gamma \times \{-1, 1\}
\]
je relacija prijelaza. Prijelaz $(q, \alpha, p, \beta, d)$ znaci:
``u stanju $q$, citajuci $\alpha$ na poziciji $m$, predi u stanje $p$,
zapisi $\beta$ na poziciju $m$, te se pomakni na poziciju $m + d$
(ali ne izlazi s trake: $0 \leq m + d < |w|$).''

OA prihvaća riječ $w$ ako iz početne konfiguracije $(q_0, 0, w)$
može doći do neke završne konfiguracije $(q_X, m, t)$.
\end{definicija}

\begin{napomena}
Kljucne razlike od PA:
\begin{itemize}
    \item Traka je ograničena na $|w|$ znakova --- nema više.
    \item Može se kretati lijevo i desno (PA samo naprijed).
    \item Može pisati po traci (modifikacija ulaza).
    \item Kao i PA: OA je \textbf{nedeterminističan}, a
      deterministični i nedeterministični OA \emph{nisu} ekvivalentni
      (za razliku od KA) --- ovo je otvoreni problem!
\end{itemize}
\end{napomena}

\begin{primjer}[OA za $L_\equiv$]
Pseudokod OA koji prepoznaje $L_\equiv = \{a^n b^n c^n\}$:
\begin{enumerate}
    \item Označi prvi $a$ i pročitaj sva $a$ dalje. Ako nema $a$, odbij.
    \item Označi prvi $b$ iza svih $a$-ova. Ako nema $b$, odbij.
    \item Označi prvi $c$ iza svih $b$-ova. Ako nema $c$, odbij.
    \item Vrati se lijevo do zadnjeg označenog $a$.
    \item Ponovi od koraka 1 s prvim neoznačenim $a$.
    \item Kad su sva $a$ označena: provjeri da su sva $b$ i $c$
      također označena (ne smije biti neoznačenih). Prihvati.
\end{enumerate}
Svaki prolaz označi točno jedno $a$, jedno $b$ i jedno $c$,
sto osigurava $|a| = |b| = |c|$.
\end{primjer}

\subsection{Odlučivost $A_\text{OA}$}

Za razliku od Turingovih strojeva, za OA uvijek postoji konačno
mnogo konfiguracija (duljina je fiksna $|w|$), pa možemo problem
prihvaćanja riješiti obilazak grafa:

Za OA $O$ i rijec $w$, konstruiramo konacni graf ciji su vrhovi
sve moguce konfiguracije duljine $|w|$:
\[
  \text{broj vrhova} = |Q| \cdot |w| \cdot |\Gamma|^{|w|}.
\]
Problem prihvaćanja svodi se na dostupnost završnih konfiguracija
iz početne --- što rješavamo DFS-om ili BFS-om.

\begin{napomena}
Graf je ogroman (eksponencijalan u $|w|$), pa algoritam nije
polinoman. Ali je konačan, pa je problem \emph{odlučiv}.
\end{napomena}

% ================================================================
\section{Poboljsanja OA}
% ================================================================

Osnovna definicija OA je praktično nezgrapna. Pokazujemo niz
ekvivalentnih poboljšanja:

\begin{center}
\renewcommand{\arraystretch}{1.5}
\begin{tabular}{ll}
\hline
\textbf{Poboljsanje} & \textbf{Ideja simulacije} \\
\hline
Nemijenjanje trake & svaki $(q, *, p, *, d)$ postaje $|\Gamma|$ prijelaza \\
Više tragova & Kartezijev produkt abeceda tragova \\
Detekcija rubova & trag s oznakama $\{\wedge, \sqcup\}$ \\
Stajanje na mjestu ($d = 0$) & lijevo pa desno, uz provjeru ruba \\
Više završnih stanja & novo stanje $q_X$, $\eps$-prijelazi iz svakog $q \in F$ \\
Konačne varijable & produkt $Q \times S$ za varijablu iz $S$ \\
Više trakova (s gl.\ nezavisnim) & $2k$ tragova, svaka glava kao oznaka \\
\hline
\end{tabular}
\end{center}

Svako od ovih poboljšanja ekvivalentno je originalnom OA u smislu
da prepoznaje istu klasu jezika. Dokazi idu metodom \textbf{simulacije}:
poboljšani stroj uvijek može sve što može obični OA (trivijalni smjer),
a obični OA simulira poboljšanog imitirajući svaki elementarni korak.

% ================================================================
\section{Kontekstne i monotone gramatike}
% ================================================================

\begin{definicija}[Kontekstna gramatika]
Pravilo $uAv \to ubv$ je \textbf{kontekstno} ako su
$u, v \in Y^*$, $A \in V$ i $b \in Y^+$. Dakle: varijabla $A$ moze
se zamijeniti s $b$ samo u ``kontekstu'' $u \cdot v$ oko nje.

\textbf{Kontekstna gramatika} (KG) je gramatika u kojoj su sva pravila
kontekstna. Jezik je \textbf{kontekstan} ako ga generira neka KG.
\end{definicija}

\begin{definicija}[Monotona gramatika]
Pravilo $u \to v$ je \textbf{monotono} ako $|u| \leq |v|$ --- desna strana
nije kraca od lijeve. Gramatika je \textbf{monotona} (MG) ako su
joj sva pravila monotona.
\end{definicija}

\begin{propozicija}
Klase $\mathrm{KG}$, $\mathrm{MG}$ i $\mathrm{OA}$ generiraju/prepoznaju
istu klasu jezika: \textbf{kontekstne jezike} Kon.
Vrijedi $\mathrm{BK}_+ \subsetneq \mathrm{Kon}$ (svi pozitivni BK-jezici
su kontekstni, ali $L_\equiv$ je kontekstan a nije BK).
\end{propozicija}

\subsection{Monotona gramatika za $L_\equiv$}

Primjer monotone gramatike koja generira $L_\equiv$:
\[
  S \to abc \mid aSBc, \quad cB \to Bc, \quad bB \to bb.
\]
Izvod $a^n b^n c^n$ koristi pravilo $S \to aSBc$ tocno $n-1$ puta,
zatim $S \to abc$, pa ``proburbla'' svaki $B$ ulijevo pokraj $c$-ova
(pravilo $cB \to Bc$), i na kraju pretvori svaki $B$ u $b$ (pravilo $bB \to bb$).

Da su pravila monotona lako se provjeriti: sve desne strane su iste
ili dulje od lijevih. Da gramatika generira točno $L_\equiv$ dokazujemo
\textbf{invarijantom}: riječ $w$ u izvodu je ``balansirana'' ako ima jednak
broj $a$-ova, $b$-ova, $B$-ova i $c$-ova zajedno gledano.

\subsection{Pretvorba MG $\to$ KG}

Pretvorba ide u dvije faze analogno pretvorbi BKG u ChNF:
\begin{enumerate}
  \item \textbf{TERM}: znakove koji se pojavljuju uz ostalo zamijenimo
        novim varijablama (i na lijevoj i na desnoj strani!).
  \item \textbf{BIN}: pravila s $n \geq 2$ simbola s lijeve strane
        $A_1 \cdots A_n \to B_1 \cdots B_m$ (gdje $m \geq n \geq 2$)
        zamijenimo s $2n$ kontekstnih pravila koja ``postepeno''
        zamjenjuju $A_i$ s $C_i$, pa $C_i$ s $B_i$.
\end{enumerate}

% ================================================================
\section{Ekvivalentnost OA i MG}
% ================================================================

\subsection{MG $\to$ OA}

Za MG $G$, OA nedeterministički ``pogađa'' izvod riječi:
koristi drugi trag za međuvriječi, primjenjuje pravila u nedeterminis-
tičnom redoslijedu, i prihvaća kad se drugi trag poklopi s ulazom.
Ključno: jer su pravila monotona, međuvriječ nikad nije dulji od
ulaza --- dakle stane na traci.

\subsection{OA $\to$ MG}

Ovaj smjer je tehnički složeniji. Koristimo \textbf{jednostavni OA} (JOA)
koji prihvaća $w$ samo u konfiguraciji $(q_X, 0, w)$ --- jedinstvenoj
završnoj konfiguraciji.

Ideja: kodiramo konfiguracije OA kao rijeci nad $Y = Q \cup \Gamma$,
jer svaka konfiguracija $(q, m, t_0 \cdots t_{n-1})$ odgovara rijeci
$t_0 \cdots t_{m-1} \, q \, t_m \cdots t_{n-1}$.

Pravila gramatike simuliraju prijelaze OA:
\begin{itemize}
  \item Za prijelaz $(q, \alpha, p, \beta, 1)$: pravilo $q\alpha \to \beta p$.
  \item Za prijelaz $(q, \alpha, p, \beta, -1)$: pravilo $\gamma q\alpha \to p\gamma\beta$
        za svaki $\gamma \in \Gamma$.
\end{itemize}

Dodatni trik: objedinjavanje znakova (``sljepljivanje'' pomocnih
simbola s pravcim znakovima) osigurava da gramatika ostane monotona
dok uklanjamo rubne oznake.

% ================================================================
\section{Zatvorenost klase Kon}
% ================================================================

\begin{propozicija}
Klasa $\mathrm{Kon}$ zatvorena je na sve skupovne i jezicne operacije:
$\cup$, $\cap$, $^c$ (komplement), $\setminus$, $\cdot$, $^+$.
\end{propozicija}

Zatvorenost na uniju, konkatenaciju i plus slijedi identicki kao za BK
(dodavanjem novog početnog simbola).

Zatvorenost na \textbf{presjek} slijedi iz ``serijskog spajanja'' JOA:
završno stanje prvog postaje početno stanje drugog.
Budući da JOA čuva traku i poziciju, drugi automat počinje točno
u stanju u kojemu je prvi zavrsio.

Zatvorenost na \textbf{komplement} je tezi:

\begin{propozicija}[Immerman--Szelepcs\'enyi, 1988]
Komplement kontekstnog jezika je kontekstan.
\end{propozicija}

Ovo je netrivi jalan teorem (dokazan tek 1988., neovisno od dva
tako sto pokaze da \emph{ne postoji} prihvacajuca grana za $w$.
autora). Ideja: nedeterministični OA koji prepoznaje komplement
koristi tehniku ``brojenja'' --- nedeterministički prihvaća $w \notin L$
tako što pokaže da \emph{ne postoji} prihvaćajuća grana za $w$.
tako sto pokaze da \emph{ne postoji} prihvacajuca grana za $w$.

% ================================================================
\section{Usporedba klasa}
% ================================================================

\begin{center}
\renewcommand{\arraystretch}{1.6}
\begin{tabular}{lllll}
\hline
& \textbf{Reg} & \textbf{BK} & \textbf{Kon} & \textbf{RE} \\
\hline
Automat & KA & PA & OA & TP \\
Gramatika & DLG/LLG & BKG & KG/MG & OG \\
  Memorija & ništa & stog & ogr.\ traka & neogr.\ traka \\
Zatvorena $\cup$ & \checkmark & \checkmark & \checkmark & \checkmark \\
Zatvorena $\cap$ & \checkmark & \texttimes & \checkmark & \checkmark \\
Zatvorena $^c$ & \checkmark & \texttimes & \checkmark & \texttimes \\
 \(A_L\) odlučiv? & \checkmark & \checkmark & \checkmark & samo poluodl. \\
 \(E_L\) odlučiv? & \checkmark & \checkmark & \texttimes & \texttimes \\
 \(\mathrm{EQ}_L\) odlučiv? & \checkmark & \texttimes & \texttimes & \texttimes \\
\hline
\end{tabular}
\end{center}

\begin{zadatak}
\begin{enumerate}[label=(\alph*), itemsep=4pt]
  \item Konstruirajte OA koji prepoznaje $\{ww^R \mid w \in \{a,b\}^*\}$.
  \item Napišite monotonu gramatiku za $\{a^n b^n c^n d^n \mid n \geq 1\}$.
  \item Dokazite da je $\{a^{2^n} \mid n \geq 0\}$ kontekstan jezik.
  \item Zasto klasa Kon nije zatvorena na Kleenevu zvijezdu?
        Koji je problem s $\eps$?
\end{enumerate}
\end{zadatak}

% ================================================================
\section*{Sazetak}
% ================================================================

\medskip\noindent
\begin{tabular}{ll}
  OA & ograniceni automat: traka duljine $|w|$, kretanje lijevo-desno \\
  JOA & jednostavni OA: jedinstven zavrsni uvjet \\
  KG & kontekstna gramatika: $uAv \to ubv$ \\
  MG & monotona gramatika: $|u| \leq |v|$ \\
  $\mathrm{OA} \equiv \mathrm{KG} \equiv \mathrm{MG}$ & svi opisuju klasu Kon \\
  I-S teorem & $\mathrm{Kon}$ zatvorena na komplement \\
\end{tabular}

\medskip\noindent
U sljedećem poglavlju prelazimo na \textbf{Turingove strojeve} ---
strojeve bez ikakvog ograničenja na traku --- i pitanja odlučivosti.
% Chapter 8 — includable by main.tex
\chapter{Turingovi strojevi i odlučivost}
\label{chap:8}

\begin{quote}
\itshape
Sve što računala mogu raditi, mogu i Turingovi strojevi --- i obrnuto.
Ovo poglavlje definira taj univerzalni model izračuna, klasificira jezike
na odlučive i neodlučive, te dokazuje temeljne rezultate neodlučivosti
pomoću dijagonalizacije i svođenja.
\end{quote}


% ================================================================
\section{Turingov stroj --- motivacija}
% ================================================================

OA je mocniji od PA, ali jos uvijek ima ogranicenje: traka ne smije
biti dulja od ulaza. Što ako trebamo više prostora?

Najjednostavnija opcija: ``otvorimo'' traku s desne strane. Umjesto
da je traka duljine $|w|$, traka je \emph{beskonačna} prema desno,
inicijalno popunjena posebnim znakom praznine $\sqcup$.

Ovaj model zovemo \textbf{Turingov stroj} (TS), i ispostavlja se da je
to upravo dovoljno za modeliranje svakog izračuna koji je ikada
opisan. Ovo je sadržaj \textbf{Church-Turingove teze}: sve što je
``algoritamski'' izračunljivo, izračunljivo je i na Turingovom stroju.

TS dolazi u tri varijante s različitim svrhama:

\begin{center}
\renewcommand{\arraystretch}{1.4}
\begin{tabular}{lll}
\hline
\textbf{Vrsta} & \textbf{Svrha} & \textbf{Ključno} \\
\hline
Prepoznavac (TP) & prepoznaje jezik & može se vrtjeti unedogled \\
Odlučitelj (TO) & odlučuje problem & uvijek stane \\
Nabrajac (TE) & generira jezik & nema ulaza, nabrajamo izlaz \\
\hline
\end{tabular}
\end{center}

% ================================================================
\section{Turingov prepoznavac (TP)}
% ================================================================

\begin{definicija}[Turingov prepoznavac]
\textbf{Turingov prepoznavac} nad $\Sig$ je sedmorka
\[
  T = (Q,\, \Sig,\, \Gamma,\, \sqcup,\, \delta,\, q_0,\, q_X),
\]
gdje je $\Sig \subsetneq \Gamma$ (radna abeceda je sira od ulazne),
$\sqcup \in \Gamma \setminus \Sig$ je znak praznine, a
\[
  \delta : (Q \setminus \{q_X\}) \times \Gamma \to Q \times \Gamma \times \{-1, 1\}
\]
je \textbf{deterministicna} funkcija prijelaza.

Konfiguracija je trojka $(q, n, t) \in Q \times \mathbb{N} \times \Gamma^\mathbb{N}_\sqcup$,
gdje je $\Gamma^\mathbb{N}_\sqcup$ skup svih traka koje su od nekog mjesta
nadalje konstanta $\sqcup$ (prebrojivo mnogo takvih).

$T$ prihvaća riječ $w$ ako izračunavanje počevši od
$(q_0, 0, w\sqcup\sqcup\cdots)$ dosegne konfiguraciju sa stanjem $q_X$.
\end{definicija}

\begin{definicija}[Rekurzivno prebrojivi jezici]
Jezik je \textbf{rekurzivno prebrojiv} (u klasi RE) ako ga prepoznaje
neki TP. Piše se $L \in \mathrm{RE}$.
\end{definicija}

Ključna razlika od OA: TP se \emph{ne mora zaustaviti} za riječi koje
ne prihvaća. Može se vrtjeti unedogled. To je problem: ne znamo
je li stroj ``razmišlja'' ili je u beskonačnoj petlji.

% ================================================================
\section{Turingov odlučitelj (TO)}
% ================================================================

\begin{definicija}[Turingov odlučitelj i rekurzivni jezici]
\textbf{Turingov odlučitelj} (TO) je TP koji uz završno stanje $q_X$
ima i stanje odbijanja $\overline{q_X}$, te se \textbf{uvijek} zaustavi ---
ili u $q_X$ (prihvaćanje) ili u $\overline{q_X}$ (odbijanje).

Jezik je \textbf{rekurzivan} (u klasi Rek) ako ga prepoznaje neki TO.
\end{definicija}

Svaki TO je TP (možemo ignorirati $\overline{q_X}$), pa $\mathrm{Rek} \subseteq \mathrm{RE}$.

\begin{propozicija}[Postov teorem]
$L \in \mathrm{Rek}$ ako i samo ako su i $L$ i $L^c$ u $\mathrm{RE}$.
\end{propozicija}

\textit{Dokaz.}
$(\Rightarrow)$: ako je $L$ rekurzivan, i $L$ i $L^c$ su rekurzivni
(Rek je zatvorena na komplement), pa a fortiori u RE.

$(\Leftarrow)$: neka TP $S$ prepoznaje $L$, a TP $T$ prepoznaje $L^c$.
Konstruiramo TO koji istovremeno simulira $S$ i $T$, korak po korak.
Jer je svaka riječ ili u $L$ ili u $L^c$, jedan od njih će sigurno stati.
Kad $S$ prihvati, TO prihvaća; kad $T$ prihvati, TO odbija. \qed

\begin{kljucno}
Postov teorem znači: jezik je odlučiv točno onda kad možemo i
\emph{potvrditi} i \emph{zanijekati} članstvo algoritmom koji uvijek stane.
Samo znati "kako naći" nije dovoljno --- treba znati i "kako ne naći".
\end{kljucno}

% ================================================================
\section{Zatvorenost klasa Rek i RE}
% ================================================================

\begin{propozicija}
Klasa $\mathrm{Rek}$ zatvorena je na sve skupovne i jezicne operacije.
Klasa $\mathrm{RE}$ zatvorena je na $\cup$ i $\cap$, ali \textbf{ne} na komplement.
\end{propozicija}

Ključna tehnika za Rek: \textbf{potpuna simulacija}. Naredba
``simuliraj rad od $T$ na sadržaju trake $t_i$'' pokreće TO na
kopiji trake, te uvijek stane jer je $T$ odlučitelj.

Primjer --- zatvorenost Rek na Kleeneovu zvijezdu:
za riječ $w$, probamo sve moguće rastave $w = w_1 \cdots w_k$ (konačno
mnogo) i za svaki rastav provjerimo je li svaki $w_i \in L$. Ako ikoji
rastav uspije, prihvati; inače odbij.

Za RE s unijom: moramo simulirati oba TP \emph{naizmjenično}, korak
po korak --- ne smijemo pokrenuti prvi do kraja jer se možda nikad
ne zaustavi. Ovo je tzv.\\ "dovetailing" tehnika.

% ================================================================
\section{Kodiranje i problemi}
% ================================================================

Da bismo mogli govoriti o TS-ovima kao ulazima TS-u, trebamo ih
\textbf{kodirati} kao riječi. Koristimo \textbf{univerzalnu abecedu}
$\Sig = \{., \mathrm{I}, (, )\}$ i sustavno kodiramo:
\[
  c(n) := \mathrm{I}^n., \quad
  c(t_1, \ldots, t_k) := (c(t_1) \cdots c(t_k)), \quad
  c(\{t_1, \ldots, t_k\}) := (c(t_1) \cdots c(t_k)).
\]
Kod stroja $T$ (ili drugog matematickog objekta) oznacavamo $\langle T \rangle$,
a par $(T, w)$ oznacavamo $\langle T, w \rangle$.

\subsection{Prirodni strojevi i problemi}

\textbf{Problem} je jezik cija konkretna kodna abeceda nije bitna
\[
  A_\mathrm{KA} := \{\langle M, w \rangle : M \text{ je KA i } w \in L(M)\}
\]
je odlučiv problem (pišemo pseudokodom koji simulira KA $M$ na $w$).

% ================================================================
\section{Odlučivi i neodlučivi problemi}
% ================================================================

\subsection{Odlučivi problemi}

\begin{propozicija}
Sljedeći problemi su odlučivi:
\begin{itemize}
  \item $A_\mathrm{KA}$, $A_\mathrm{RI}$, $A_\mathrm{NKA}$ --- prihvaćanje za regularne jezike
  \item $E_\mathrm{KA}$ --- praznost KA-jezika
  \item $\mathrm{EQ}_\mathrm{KA}$ --- jednakost dva KA-jezika
  \item $A_\mathrm{BKG}$, $A_\mathrm{BK}$ --- prihvaćanje za BK-jezike (CYK!)
  \item $E_\mathrm{BKG}$, $E_\mathrm{BK}$ --- praznost BK-jezika
  \item $A_\mathrm{OA}$, $A_\mathrm{Kon}$ --- prihvaćanje za kontekstne jezike
\end{itemize}
\end{propozicija}

Svi ovi problemi imaju konkretne algoritme koje smo vec opisali
(pseudokodom ili formalnom konstrukcijom).

\subsection{Neodlučivi problemi --- hijerarhija}

\begin{center}
\renewcommand{\arraystretch}{1.5}
\begin{tabular}{lll}
\hline
\textbf{Problem} & \textbf{Klasa} & \textbf{Zašto} \\
\hline
$A_\mathrm{TP}$, $A_\mathrm{RE}$ & $\mathrm{RE} \setminus \mathrm{Rek}$ & TP može sebe simulirati \\
$\overline{A_\mathrm{TP}}$ & nije u RE & Postov teorem + $A_\mathrm{TP} \in \mathrm{RE} \setminus \mathrm{Rek}$ \\
\mathrm{EQ}_\mathrm{BK} & neodlučivo & svođenja od $\overline{A_\mathrm{TP}}$ \\
E_\mathrm{Kon} & neodlučivo & svođenja od $A_\mathrm{TP}$ \\
\hline
\end{tabular}
\end{center}

% ================================================================
\section{Svođenje (redukcija)}
% ================================================================

\begin{definicija}[Svođenje]
Kažemo da se $L$ \textbf{svodi na} $M$ ($L \leq M$) ako postoji
totalna Turing-izračunljiva funkcija $\varphi : \Sig^* \to \Sig^*$
takva da za svaki $w$:
\[
  w \in L \iff \varphi(w) \in M.
\]
\end{definicija}

\begin{propozicija}
Neka $L \leq M$.
\begin{itemize}
  \item Ako je $M$ (polu)odlučiv, tada je i $L$ (polu)odlučiv.
  \item Ako je $L$ neodlučiv, tada je i $M$ neodlučiv.
\end{itemize}
\end{propozicija}

Drugu točku koristimo za dokazivanje neodlučivosti:
Pokažemo $\mathrm{POZNATO\_NEODLUCIVO} \leq \mathrm{NOVI\_PROBLEM}$,
Čime zaključujemo da je i novi problem neodlučiv.

\begin{primjer}[Svođenje za $E_\mathrm{KA}$]
Problem jednakosti $\mathrm{EQ}_\mathrm{KA}$ svodi se na $E_\mathrm{KA}$:
\[
  \langle M_1, M_2 \rangle \mapsto \langle M_\Delta \rangle,
  \quad \text{gdje } L(M_\Delta) = L(M_1) \triangle L(M_2).
\]
Dakle $\langle M_1, M_2 \rangle \in \mathrm{EQ}_\mathrm{KA}
\iff \langle M_\Delta \rangle \in E_\mathrm{KA}$.
Konstruiranje $M_\Delta$ je mehaničko (Kartezijev produkt),
pa je $\varphi$ Turing-izračunljiva. Kako je $E_\mathrm{KA}$ odlučiv,
slijedi da je i $\mathrm{EQ}_\mathrm{KA}$ odlučiv. \checkmark
\end{primjer}

% ================================================================
\section{Dijagonalizacija: Russell i $A_\mathrm{TP}$}
% ================================================================

Kako dokazati da neki jezik \emph{nije} u RE? Lema o pumpanju
ne postoji za ovu klasu. Koristimo \textbf{dijagonalizaciju} --- trik
koji potjece od Cantora (prebrojiivost skupova) i Russella (paradoks).

\begin{propozicija}
Jezik
\[
  \mathrm{Russell} := \{\langle T \rangle : T \text{ je TP i } \langle T \rangle \notin L(T)\}
\]
nije rekurzivno prebrojiv.
\end{propozicija}

\textit{Dokaz.}
Pretpostavimo da je Russell u RE, i neka TP $R$ prepoznaje Russell.
Tada:
\[
  \langle R \rangle \in \mathrm{Russell}
  \iff \langle R \rangle \notin L(R) = \mathrm{Russell}.
\]
Kontradikcija. \qed

Ovo je direktan analog Russellovog paradoksa: ``skup svih skupova
koji ne sadrze sami sebe'' vodi do kontradikcije.

\begin{propozicija}
Jezik $A_\mathrm{TP} = \{\langle T, w \rangle : T \text{ prihvaća } w\}$
je u $\mathrm{RE} \setminus \mathrm{Rek}$.
\end{propozicija}

\textit{Skica dokaza.}
$A_\mathrm{TP} \in \mathrm{RE}$: konstruiramo \textbf{univerzalni Turingov stroj}
(UTS) koji simulira rad TP $T$ na ulazu $w$ (dan kao kod).
UTS je sam po sebi TP, pa je $A_\mathrm{TP} \in \mathrm{RE}$.

$A_\mathrm{TP} \notin \mathrm{Rek}$: dokazemo da $\overline{A_\mathrm{TP}}$ nije u RE.
Koristimo Russell $\leq \overline{A_\mathrm{TP}}$: svaki kod $\langle T \rangle$
preslikamo u $\langle T, \langle T \rangle \rangle$, i vrijedi:
\[
  \langle T \rangle \in \mathrm{Russell}
  \iff \langle T \rangle \notin L(T)
  \iff \langle T, \langle T \rangle \rangle \in \overline{A_\mathrm{TP}}.
\]
Dakle Russell $\leq \overline{A_\mathrm{TP}}$. Budući da Russell nije u RE,
ni $\overline{A_\mathrm{TP}}$ nije u RE. Po Postovom teoremu,
$A_\mathrm{TP} \notin \mathrm{Rek}$. \qed

% ================================================================
\section{Problem zaustavljanja i kompletna slika}
% ================================================================

\begin{propozicija}
Jezik $H = \{\langle T, w \rangle : \text{izračunavanje } T \text{ na } w \text{ stane}\}$
(``halting problem'') ekvivalentan je $A_\mathrm{TP}$ u smislu svedivosti:
$H \leq A_\mathrm{TP}$ i $A_\mathrm{TP} \leq H$.
\end{propozicija}

Ova ekvivalencija nije iznenadujuca: TP koji smo definirali stane
točno kada prihvate, pa "stajanje" i "prihvaćanje" su isti uvjet.

\subsection{Kompletna Chomskyjeva hijerarhija}

\begin{center}
\renewcommand{\arraystretch}{1.5}
\begin{tabular}{lllll}
\hline
\textbf{Klasa} & \textbf{Automat} & \textbf{Gram.} & \textbf{$A_L$} & \textbf{$E_L$} \\
\hline
Reg & KA & DLG & odlučiv & odlučiv \\
BK & PA & BKG & odlučiv & odlučiv \\
Kon & OA & KG/MG & odlučiv & neodlučiv \\
Rek & TO & --- & odlučiv & neodlučiv \\
RE & TP & OG & poluodlučiv & neodlučiv \\
\hline
\end{tabular}
\end{center}

Prave inkluzije potvrdjuju konkretni jezici:
\begin{align*}
  L_= &= \{a^n b^n\} \in \mathrm{BK} \setminus \mathrm{Reg} \\
  L_\equiv &= \{a^n b^n c^n\} \in \mathrm{Kon} \setminus \mathrm{BK} \\
  \mathrm{EQ}_\mathrm{RI_2} &\in \mathrm{Rek} \setminus \mathrm{Kon} \\
  A_\mathrm{TP} &\in \mathrm{RE} \setminus \mathrm{Rek} \\
  \overline{A_\mathrm{TP}} &\notin \mathrm{RE}
\end{align*}

% ================================================================
\section{Nedeterministicki Turingov stroj (NTP)}
% ================================================================

\begin{definicija}
NTP ima relaciju prijelaza umjesto funkcije. Prihvaća riječ
ako \emph{postoji barem jedna} grana izračunavanja koja stane.
\end{definicija}

\begin{propozicija}
NTP i TP su ekvivalentni: svaki NTP ima ekvivalentni TP.
\end{propozicija}

Dokaz: TP simulira NTP koristeci \textbf{tri trake}:
\begin{enumerate}
  \item Kopija ulaza (nepromjenjiva).
  \item Radna traka (simulacija jedne grane).
  \item Treća traka: kôd trenutne grane kao niz odabranih prijelaza
        iz $\Delta^*$ (leksikografskim redoslijedom isprobavamo sve grane).
\end{enumerate}
TP prolazi kroz sve moguce grane BFS-om: svaka iteracija produlji
granu za jedan korak. Čim neka grana prihvati, TP prihvaća.

\begin{napomena}
Za odlučitelje (NTO) analogna ekvivalencija vrijedi, ali je dokaz
teži --- ne može koristiti istu trotračnu simulaciju jer ne znamo
znači li $\overline{q_X}$ konačno odbijanje ili samo lokalni neuspjeh.
\end{napomena}

% ================================================================
\section{Turingovi nabrajaci (TE)}
% ================================================================

\begin{definicija}[Turingov nabrajac]
TE nema ulaza. Ima radnu traku i \textbf{izlaznu traku}, te posebno
\textbf{izlazno stanje} $q_\leftarrow$. Rad TE je uvijek beskonačni niz
konfiguracija. TE \textbf{izvodi} riječ $w$ ako ikad dosegne stanje
$q_\leftarrow$ s $w$ na izlaznoj traci.

$L(E) := \{w : E \text{ izvodi } w\}$.
\end{definicija}

\begin{propozicija}
$L \in \mathrm{RE}$ ako i samo ako postoji TE koji generira $L$.
\end{propozicija}

Smjer $(\Leftarrow)$: TP koji prepoznaje $L(E)$ simulira $E$ korak po
korak i uspoređuje izlaz s ulaznom riječi $w$. Čim se poklope,
prihvaća.

Smjer $(\Rightarrow)$: TE za $L$ (prepoznajan TP $T$) nabrajanjem:
za svaki $n$ i svaki moguc ulaz duljine $\leq n$, simulira $T$ na tom
ulazu točno $n$ koraka. Ako $T$ prihvati, izvede tu riječ.
Ovim ``dovetailing'' postupkom svaka riječ iz $L$ će biti izvedena
u konačno mnogo koraka.

% ================================================================
\section*{Sazetak}
% ================================================================

\medskip\noindent
\begin{tabular}{ll}
  TP & Turingov prepoznavač: stane ili prihvaća (nema garantije zaustavljanja) \\
  TO & Turingov odlučitelj: uvijek stane, prihvaća ili odbija \\
  RE & jezici koje prepoznaje neki TP \\
  Rek & jezici koje prepoznaje neki TO \\
  $A_\mathrm{TP}$ & poluodlučiv ali neodlučiv (svođenje + dijagonalizacija) \\
  Postov teorem & $L \in \mathrm{Rek} \iff L, L^c \in \mathrm{RE}$ \\
  Svođenje & $L \leq M$: odlučivost "prenosi se gore" \\
\end{tabular}

\medskip\noindent
\begin{center}
\Large\itshape
$\mathrm{Reg} \subsetneq \mathrm{BK} \subsetneq \mathrm{Kon} \subsetneq \mathrm{Rek} \subsetneq \mathrm{RE} \subsetneq \text{svi jezici}$
\end{center}

\end{document}
